% Generated by the blueprint tool (blueprint-pdf/1) from a blueprint snapshot; edit the
% Markdown sources, not this file.  Compile with `latexmk -pdf` or by running
% `pdflatex` twice (the second run settles the contents and the page references).
% Images, if any, are read from assets/ next to this file.
\documentclass[11pt]{article}
\usepackage[T1]{fontenc}
\usepackage[utf8]{inputenc}
\usepackage{lmodern}
\usepackage{textcomp}
\usepackage[a4paper,margin=2.5cm,headheight=14pt]{geometry}
\usepackage{amsmath,amssymb,amsthm}
\usepackage{amscd}  % KaTeX's CD diagrams; part of amsmath, so always there
\usepackage{mathtools}
\usepackage{bm}
\usepackage{graphicx}
\usepackage{xcolor}
\usepackage{array,booktabs,tabularx}
\usepackage{fancyvrb}
\usepackage[normalem]{ulem}
\usepackage{microtype}
% Optional, and only with their fonts: a small TeX installation has the
% packages without the fonts, which fails at the first formula.
\IfFileExists{mathrsfs.sty}{\IfFileExists{rsfs10.tfm}{\usepackage{mathrsfs}}{}}{}
\IfFileExists{stmaryrd.sty}{\IfFileExists{stmary10.tfm}{\usepackage{stmaryrd}}{}}{}
\IfFileExists{cancel.sty}{\usepackage{cancel}}{}
\usepackage[hyphens]{url}
\usepackage[pdfusetitle,bookmarksnumbered=false]{hyperref}
\usepackage{bookmark}
\definecolor{bplink}{HTML}{1F4E9C}
\definecolor{bpmuted}{HTML}{5F676E}
\hypersetup{colorlinks=true,linkcolor=bplink,urlcolor=bplink,citecolor=bplink}
\setlength{\parindent}{0pt}
\setlength{\parskip}{0.55em plus 0.2em minus 0.1em}
% Lists as tight as the page's.
\makeatletter
\def\@listi{\leftmargin\leftmargini\parsep 0pt\topsep 0.3em plus 0.1em\itemsep 0.1em plus 0.1em}
\let\@listI\@listi
\def\@listii{\leftmargin\leftmarginii\labelwidth\leftmarginii\advance\labelwidth-\labelsep\parsep 0pt\topsep 0.1em\itemsep 0.1em}
% An ordered list that does not start at 1.
\newcommand{\bpstartat}[1]{\setcounter{\@enumctr}{\numexpr#1-1\relax}}
\makeatother
\setcounter{tocdepth}{3}
\emergencystretch=3em
\allowdisplaybreaks

% A numbered statement, set as a paper sets one: "Theorem 2.1 (Title)."
% #1 is the lead ("Theorem~2.1"), #2 what follows it in roman.
\newenvironment{bpstatement}[2]{\par\addvspace{0.6\baselineskip}\noindent\textbf{#1}#2\textbf{.}\enspace\ignorespaces}
  {\par\addvspace{0.4\baselineskip}}
\newenvironment{bpstatementit}[2]{\begin{bpstatement}{#1}{#2}\itshape\ignorespaces}{\end{bpstatement}}
% A step: a relation with prose of its own, read under the object it starts from.
\newenvironment{bpstep}[1]{\par\addvspace{0.3\baselineskip}\begin{list}{}{\leftmargin=1.6em\rightmargin=0pt\topsep=0pt}\item\relax
  {\small\color{bpmuted}#1}\par\nopagebreak}{\end{list}\par}
% Notes about where an object is written out.
\newenvironment{bpnote}{\par\small\color{bpmuted}\itshape}{\par}
\newcommand{\bpomitted}[1]{\par{\small\color{bpmuted}\itshape #1\par}}
\newcommand{\bpbodyheading}[1]{\par\addvspace{0.4\baselineskip}\noindent\textbf{#1}\par\nopagebreak}
\newcommand{\bprule}{\par\medskip\noindent\hfil\rule{0.3\linewidth}{0.4pt}\hfil\par\medskip}
\newcommand{\bpboundary}[1]{\par\nopagebreak{\small\color{bpmuted}#1\par}\nopagebreak}
% A reference to something this export leaves out, and one that resolves nowhere.
\newcommand{\bpout}[1]{#1\textsuperscript{\hyperlink{bp-outside}{\ensuremath{\ast}}}}
\newcommand{\bpbroken}[1]{\textcolor{red!70!black}{[#1?]}}
\newcommand{\bpslug}[1]{\texttt{#1}}
% A bibliography entry: label, anchor, entry; the label hangs in the margin.
\newcommand{\bpbibitem}[3]{\par\setbox0\hbox{[#1]\enspace}\ifdim\wd0<5em\setbox0\hbox to 5em{[#1]\hfil}\fi
  \hangindent=\wd0 \hangafter=1 \noindent\box0\phantomsection\label{#2}#3\par\smallskip}
% A character the export knows no way to set.
\newcommand{\bpglyph}[1]{\fbox{\scriptsize #1}}
% Images, at their own size unless they are wider than the line.
\newsavebox{\bpimagebox}
\newcommand{\bpimage}[1]{\sbox{\bpimagebox}{\includegraphics{#1}}%
  \ifdim\wd\bpimagebox>\linewidth\resizebox{\linewidth}{!}{\usebox{\bpimagebox}}\else\usebox{\bpimagebox}\fi}
\newcommand{\bpmissingimage}[1]{\fbox{\parbox{0.8\linewidth}{\centering\small\color{bpmuted}Image not included: #1}}}
\DefineVerbatimEnvironment{bpverbatim}{Verbatim}{fontsize=\small,xleftmargin=1em}
\renewcommand{\qedsymbol}{\ensuremath{\blacksquare}}
\hypersetup{pdftitle={Weil Conjectures — 4 Proper base change},pdfsubject={Chapter 4 Proper base change; all levels},pdfcreator={blueprint}}
% KaTeX commands the website knows
\providecommand{\R}{\mathbb{R}}\providecommand{\N}{\mathbb{N}}\providecommand{\Z}{\mathbb{Z}}
\providecommand{\Q}{\mathbb{Q}}\providecommand{\C}{\mathbb{C}}
\providecommand{\Reals}{\mathbb{R}}\providecommand{\reals}{\mathbb{R}}\providecommand{\Complex}{\mathbb{C}}
\providecommand{\cnums}{\mathbb{C}}\providecommand{\natnums}{\mathbb{N}}\providecommand{\Bbb}[1]{\mathbb{#1}}
\providecommand{\bold}[1]{\mathbf{#1}}\providecommand{\frak}[1]{\mathfrak{#1}}\providecommand{\mathbbm}[1]{\mathbb{#1}}
\providecommand{\lang}{\langle}\providecommand{\rang}{\rangle}\providecommand{\bra}[1]{\langle #1\rvert}
\providecommand{\ket}[1]{\lvert #1\rangle}\providecommand{\braket}[1]{\langle #1\rangle}
\providecommand{\set}[1]{\{#1\}}\providecommand{\Set}[1]{\left\{#1\right\}}
\providecommand{\infin}{\infty}\providecommand{\isin}{\in}\providecommand{\plusmn}{\pm}\providecommand{\sub}{\subset}
\providecommand{\sube}{\subseteq}\providecommand{\supe}{\supseteq}\providecommand{\exist}{\exists}
\providecommand{\bull}{\bullet}\providecommand{\larr}{\leftarrow}\providecommand{\rarr}{\rightarrow}
\providecommand{\lrarr}{\leftrightarrow}\providecommand{\Larr}{\Leftarrow}\providecommand{\Rarr}{\Rightarrow}
\providecommand{\Lrarr}{\Leftrightarrow}\providecommand{\uarr}{\uparrow}\providecommand{\darr}{\downarrow}
\providecommand{\harr}{\leftrightarrow}\providecommand{\lArr}{\Leftarrow}\providecommand{\rArr}{\Rightarrow}
\providecommand{\hArr}{\Leftrightarrow}\providecommand{\alef}{\aleph}\providecommand{\alefsym}{\aleph}
\providecommand{\real}{\Re}\providecommand{\image}{\Im}\providecommand{\weierp}{\wp}
\providecommand{\Alpha}{\mathrm{A}}\providecommand{\Beta}{\mathrm{B}}\providecommand{\Epsilon}{\mathrm{E}}
\providecommand{\Zeta}{\mathrm{Z}}\providecommand{\Eta}{\mathrm{H}}\providecommand{\Iota}{\mathrm{I}}
\providecommand{\Kappa}{\mathrm{K}}\providecommand{\Mu}{\mathrm{M}}\providecommand{\Nu}{\mathrm{N}}
\providecommand{\Omicron}{\mathrm{O}}\providecommand{\Rho}{\mathrm{P}}\providecommand{\Tau}{\mathrm{T}}
\providecommand{\Chi}{\mathrm{X}}\providecommand{\omicron}{o}\providecommand{\gt}{>}\providecommand{\lt}{<}
\providecommand{\And}{\&}\providecommand{\hdashline}{\hline}\providecommand{\operatornamewithlimits}[1]{\operatorname*{#1}}
\providecommand{\llbracket}{[\![}\providecommand{\rrbracket}{]\!]}\providecommand{\cancel}[1]{#1}
\providecommand{\mathscr}{\mathcal}
% the characters beyond ASCII this document uses
\DeclareUnicodeCharacter{2160}{\textup{I}}
\DeclareUnicodeCharacter{2161}{\textup{II}}
% the project's macros ([katex.macros] in blueprint.toml)
\providecommand{\A}{}\renewcommand{\A}{\mathbf A}
\providecommand{\AEn}{}\renewcommand{\AEn}{A_{E,n}}
\providecommand{\Aaf}{}\renewcommand{\Aaf}{\mathbf A}
\providecommand{\Ab}{}\renewcommand{\Ab}{\mathsf{Ab}}
\providecommand{\AffSch}{}\renewcommand{\AffSch}{\operatorname{AffSch}}
\providecommand{\Aone}{}\renewcommand{\Aone}{\mathbf A^1}
\providecommand{\Bl}{}\renewcommand{\Bl}{\operatorname{Bl}}
\providecommand{\Br}{}\renewcommand{\Br}{\mathrm{Br}}
\providecommand{\C}{}\renewcommand{\C}{\mathbb C}
\providecommand{\Cech}{}\renewcommand{\Cech}{\check{H}}
\providecommand{\Coext}{}\renewcommand{\Coext}{\operatorname{Coext}}
\providecommand{\Cone}{}\renewcommand{\Cone}{\operatorname{Cone}}
\providecommand{\Cons}{}\renewcommand{\Cons}{\operatorname{Cons}}
\providecommand{\Cref}{}\renewcommand{\Cref}[1]{\text{#1}}
\providecommand{\D}{}\renewcommand{\D}{D}
\providecommand{\DLAFT}{}\renewcommand{\DLAFT}{\operatorname{FT}}
\providecommand{\DLAFr}{}\renewcommand{\DLAFr}{\operatorname{Fr}}
\providecommand{\DLAQ}{}\renewcommand{\DLAQ}{{\overline{\mathbf Q}_\ell}}
\providecommand{\DLRQ}{}\renewcommand{\DLRQ}{{\overline{\mathbf Q}_\ell}}
\providecommand{\Dbc}{}\renewcommand{\Dbc}{D^b_c}
\providecommand{\Dc}{}\renewcommand{\Dc}{D_{\mathrm{cons}}}
\providecommand{\Dctf}{}\renewcommand{\Dctf}{D^b_{ctf}}
\providecommand{\Dic}{}\renewcommand{\Dic}{D_{\IndCons}}
\providecommand{\Dplus}{}\renewcommand{\Dplus}{D^{+}}
\providecommand{\End}{}\renewcommand{\End}{\operatorname{End}}
\providecommand{\Et}{}\renewcommand{\Et}{\mathrm{\acute{e}t}}
\providecommand{\Ex}{}\renewcommand{\Ex}{\mathrm{Ex}}
\providecommand{\Ext}{}\renewcommand{\Ext}{\operatorname{Ext}}
\providecommand{\F}{}\renewcommand{\F}{\mathbb F}
\providecommand{\FT}{}\renewcommand{\FT}{\operatorname{FT}}
\providecommand{\Fbar}{}\renewcommand{\Fbar}{{\overline{\mathbf F}_q}}
\providecommand{\Fell}{}\renewcommand{\Fell}{\mathbf F_\ell}
\providecommand{\Fib}{}\renewcommand{\Fib}{\operatorname{Fib}}
\providecommand{\Fix}{}\renewcommand{\Fix}{\operatorname{Fix}}
\providecommand{\Fp}{}\renewcommand{\Fp}{{\mathbf F_{p}}}
\providecommand{\Fq}{}\renewcommand{\Fq}{{\mathbf F_q}}
\providecommand{\Fqn}{}\renewcommand{\Fqn}{{\mathbf F_{q^{n}}}}
\providecommand{\Frob}{}\renewcommand{\Frob}{\operatorname{Frob}}
\providecommand{\GL}{}\renewcommand{\GL}{\operatorname{GL}}
\providecommand{\Gal}{}\renewcommand{\Gal}{\operatorname{Gal}}
\providecommand{\Hens}{}\renewcommand{\Hens}{\operatorname{Hens}}
\providecommand{\Hilb}{}\renewcommand{\Hilb}{\operatorname{Hilb}}
\providecommand{\Hom}{}\renewcommand{\Hom}{\operatorname{Hom}}
\providecommand{\IndCons}{}\renewcommand{\IndCons}{\operatorname{IndCons}}
\providecommand{\Kar}{}\renewcommand{\Kar}{\operatorname{Kar}}
\providecommand{\LT}{}\renewcommand{\LT}{\operatorname{LT}}
\providecommand{\Lie}{}\renewcommand{\Lie}{\operatorname{Lie}}
\providecommand{\Loc}{}\renewcommand{\Loc}{\operatorname{Loc}}
\providecommand{\Mod}{}\renewcommand{\Mod}{\mathsf{Mod}}
\providecommand{\N}{}\renewcommand{\N}{\mathbb{N}}
\providecommand{\Needspace}{}\renewcommand{\Needspace}{\par}
\providecommand{\O}{}\renewcommand{\O}{{\mathcal{O}}}
\providecommand{\Ob}{}\renewcommand{\Ob}{\mathrm{Ob}}
\providecommand{\Oe}{}\renewcommand{\Oe}{\mathcal O_E}
\providecommand{\PP}{}\renewcommand{\PP}{\mathbf{P}}
\providecommand{\Perf}{}\renewcommand{\Perf}{\operatorname{Perf}}
\providecommand{\Pic}{}\renewcommand{\Pic}{\mathrm{Pic}}
\providecommand{\Pm}{}\renewcommand{\Pm}{\mathbf P^{m}}
\providecommand{\Pone}{}\renewcommand{\Pone}{\mathbf P^{1}}
\providecommand{\Pqn}{}\renewcommand{\Pqn}{\mathbf P^{n}}
\providecommand{\ProFin}{}\renewcommand{\ProFin}{\operatorname{ProFin}}
\providecommand{\Proet}{}\renewcommand{\Proet}{\mathrm{pro\acute{e}t}}
\providecommand{\Q}{}\renewcommand{\Q}{\mathbb Q}
\providecommand{\Qb}{}\renewcommand{\Qb}{\overline{\mathbf Q}_{\ell}}
\providecommand{\Qbar}{}\renewcommand{\Qbar}{\overline{\mathbf Q}}
\providecommand{\QbarX}{}\renewcommand{\QbarX}{\overline{\mathbf Q}_{\ell,X}}
\providecommand{\Qbarl}{}\renewcommand{\Qbarl}{{\overline{\mathbf Q}_{\ell}}}
\providecommand{\Qell}{}\renewcommand{\Qell}{{\mathbf{Q}_{\ell}}}
\providecommand{\Ql}{}\renewcommand{\Ql}{\mathbf Q_{\ell}}
\providecommand{\Qlb}{}\renewcommand{\Qlb}{{\overline{\mathbf Q}_\ell}}
\providecommand{\Qlbar}{}\renewcommand{\Qlbar}{\overline{\mathbf Q}_\ell}
\providecommand{\R}{}\renewcommand{\R}{\mathbb{R}}
\providecommand{\RG}{}\renewcommand{\RG}{\operatorname{\mathrm{R}\Gamma}}
\providecommand{\RGamma}{}\renewcommand{\RGamma}{R\Gamma}
\providecommand{\RGammac}{}\renewcommand{\RGammac}{R\Gamma_c}
\providecommand{\RHom}{}\renewcommand{\RHom}{R\mathcal{H}om}
\providecommand{\RHomX}{}\renewcommand{\RHomX}{R\mathcal H om_X}
\providecommand{\Rat}{}\renewcommand{\Rat}{\operatorname{IsRat}}
\providecommand{\Res}{}\renewcommand{\Res}{\operatorname{Res}}
\providecommand{\Rlim}{}\renewcommand{\Rlim}{R\!\varprojlim}
\providecommand{\SIm}{}\renewcommand{\SIm}{\operatorname{SIm}}
\providecommand{\Sch}{}\renewcommand{\Sch}{\mathsf{Sch}}
\providecommand{\Sh}{}\renewcommand{\Sh}{\operatorname{Sh}}
\providecommand{\Span}{}\renewcommand{\Span}{\operatorname{span}}
\providecommand{\Spec}{}\renewcommand{\Spec}{\operatorname{Spec}}
\providecommand{\Tr}{}\renewcommand{\Tr}{\operatorname{Tr}}
\providecommand{\Trsm}{}\renewcommand{\Trsm}{\operatorname{Tr}^{\mathrm{sm}}}
\providecommand{\Wt}{}\renewcommand{\Wt}{\operatorname{Wt}}
\providecommand{\Z}{}\renewcommand{\Z}{\mathbb{Z}}
\providecommand{\Zar}{}\renewcommand{\Zar}{\mathrm{zar}}
\providecommand{\Zell}{}\renewcommand{\Zell}{{\mathbf{Z}_{\ell}}}
\providecommand{\Zl}{}\renewcommand{\Zl}{{\mathbf Z_\ell}}
\providecommand{\affet}{}\renewcommand{\affet}[1]{{#1}^{\mathrm{aff}}_{\mathrm{ét}}}
\providecommand{\affproet}{}\renewcommand{\affproet}[1]{{#1}^{\mathrm{aff}}_{\mathrm{proét}}}
\providecommand{\boxt}{}\renewcommand{\boxt}{\boxtimes_A^{\mathrm L}}
\providecommand{\cF}{}\renewcommand{\cF}{\mathcal{F}}
\providecommand{\cHom}{}\renewcommand{\cHom}{R\mathcal{H}\!om}
\providecommand{\cI}{}\renewcommand{\cI}{\mathcal{I}}
\providecommand{\cL}{}\renewcommand{\cL}{\mathcal L}
\providecommand{\cM}{}\renewcommand{\cM}{\mathcal M}
\providecommand{\cN}{}\renewcommand{\cN}{\mathcal N}
\providecommand{\cO}{}\renewcommand{\cO}{\mathcal O}
\providecommand{\cT}{}\renewcommand{\cT}{\mathcal T}
\providecommand{\checked}{}\renewcommand{\checked}{\color{red}\aftergroup\restorecolor}
\providecommand{\coeff}{}\renewcommand{\coeff}{\operatorname{coeff}}
\providecommand{\coker}{}\renewcommand{\coker}{\operatorname{coker}}
\providecommand{\colim}{}\renewcommand{\colim}{\operatorname{colim}}
\providecommand{\constructedby}{}\renewcommand{\constructedby}[1]{\par\smallskip\noindent\textcolor{purple!60!black}{\textsf{Constructed by:} }
  \footnotesize\sloppy\nolinkurl{#1}\par}
\providecommand{\deepresult}{}\renewcommand{\deepresult}[1]{\par\smallskip\noindent\textcolor{orange!45!black}{
  \textsf{Blueprint role: named deep theorem node; see the cluster proof architecture.}}
  \par\noindent\textcolor{orange!45!black}{\textsf{Mathematical source:} }#1\par}
\providecommand{\dt}{}\renewcommand{\dt}{\mathbin{\otimes^{\mathbf L}}}
\providecommand{\dw}{}\renewcommand{\dw}{\operatorname{dw}}
\providecommand{\et}{}\renewcommand{\et}[1]{{#1}_{\mathrm{ét}}}
\providecommand{\ev}{}\renewcommand{\ev}{\operatorname{ev}}
\providecommand{\fib}{}\renewcommand{\fib}{\operatorname{fib}}
\providecommand{\formalizationnote}{}\renewcommand{\formalizationnote}[1]{\begin{quote}\small\raggedright\textbf{Formalization note.} #1\end{quote}}
\providecommand{\fpqc}{}\renewcommand{\fpqc}{X_{\mathrm{fpqc}}}
\providecommand{\gammacont}{}\renewcommand{\gammacont}{\Gamma_{\mathrm{cont}}}
\providecommand{\gm}{}\renewcommand{\gm}{\mathbf G_m}
\providecommand{\hocolim}{}\renewcommand{\hocolim}{\operatorname{hocolim}}
\providecommand{\holim}{}\renewcommand{\holim}{\operatorname{holim}}
\providecommand{\hrB}{}\renewcommand{\hrB}{\@ifnextchar]{\hrEnd}{\hrStep}}
\providecommand{\hrEnd}{}\renewcommand{\hrEnd}{\@firstoftwo{}}
\providecommand{\hrStep}{}\renewcommand{\hrStep}{\@firstoftwo{\hrB}}
\providecommand{\id}{}\renewcommand{\id}{\operatorname{id}}
\providecommand{\kunFl}{}\renewcommand{\kunFl}{\mathbf F_\ell}
\providecommand{\kunZn}{}\renewcommand{\kunZn}{\mathbf Z/n\mathbf Z}
\providecommand{\kuncd}{}\renewcommand{\kuncd}{\operatorname{cd}}
\providecommand{\kundim}{}\renewcommand{\kundim}{\dim}
\providecommand{\kunotimes}{}\renewcommand{\kunotimes}[1]{\otimes^{\mathbf L}_{#1}}
\providecommand{\kunsupp}{}\renewcommand{\kunsupp}{\operatorname{Supp}}
\providecommand{\lotimes}{}\renewcommand{\lotimes}{\otimes^{\mathrm L}}
\providecommand{\m}{}\renewcommand{\m}{{\mathfrak{m}}}
\providecommand{\n}{}\renewcommand{\n}{{\mathfrak{n}}}
\providecommand{\nagataGap}{}\renewcommand{\nagataGap}[1]{\par\noindent\textbf{Formalization gap.} #1\par}
\providecommand{\openobligation}{}\renewcommand{\openobligation}[1]{\par\smallskip\noindent\textcolor{red!65!black}{\textsf{Open obligation:} }
  \footnotesize\sloppy\nolinkurl{#1}\par}
\providecommand{\p}{}\renewcommand{\p}{{\mathfrak{p}}}
\providecommand{\pr}{}\renewcommand{\pr}{\operatorname{pr}}
\providecommand{\preshv}{}\renewcommand{\preshv}[1]{\mathrm{PreShv}(#1)}
\providecommand{\pro}{}\renewcommand{\pro}[1]{\mathrm{Pro}(#1)}
\providecommand{\proet}{}\renewcommand{\proet}{X_{\mathrm{proét}}}
\providecommand{\q}{}\renewcommand{\q}{{\mathfrak{q}}}
\providecommand{\rad}{}\renewcommand{\rad}{\operatorname{rad}}
\providecommand{\rank}{}\renewcommand{\rank}{\operatorname{rank}}
\providecommand{\red}{}\renewcommand{\red}{\color{red}}
\providecommand{\restorecolor}{}\renewcommand{\restorecolor}{\color{black}}
\providecommand{\retiredobligation}{}\renewcommand{\retiredobligation}[1]{\par\smallskip\noindent\textcolor{red!75!black}{\textsf{Retired invalid route:} }
  \footnotesize\sloppy\nolinkurl{#1}\par}
\providecommand{\rk}{}\renewcommand{\rk}{\operatorname{rk}}
\providecommand{\sep}{}\renewcommand{\sep}{\textnormal{sep}}
\providecommand{\shv}{}\renewcommand{\shv}[1]{\mathrm{Shv}(#1)}
\providecommand{\spec}{}\renewcommand{\spec}[1]{\mathrm{Spec}(#1)}
\providecommand{\stacks}{}\renewcommand{\stacks}[1]{\href{https://stacks.math.columbia.edu/tag/#1}{\textnormal{Stacks, Tag~#1}}}
\providecommand{\stacksproject}{}\renewcommand{\stacksproject}[1]{{\cite[\href{https://stacks.math.columbia.edu/tag/#1}{Tag #1}]{stacks-project}}}
\providecommand{\statusnote}{}\renewcommand{\statusnote}[1]{\begin{quote}\small\raggedright\textbf{Status.} #1\end{quote}}
\providecommand{\tildeideal}{}\renewcommand{\tildeideal}[2]{\widetilde{(#1, #2)}}
\providecommand{\tildering}{}\renewcommand{\tildering}[3]{{#1}_{\widetilde{(#2, #3)}}}
\providecommand{\tr}{}\renewcommand{\tr}{\operatorname{tr}}
\providecommand{\wt}{}\renewcommand{\wt}{\operatorname{wt}}

\begin{document}
\pagestyle{plain}
\begin{center}
{\LARGE\bfseries Weil Conjectures\par}
\vspace{0.7em}
{\large Chapter 4: Proper base change\par}
\vspace{0.5em}
{\small\color{bpmuted}Snapshot 2026-09-26, revision 8e4cb88e2a\par}
\end{center}
\vspace{0.5em}
{\small\tableofcontents}
\bigskip\hrule\bigskip

% object: chap-Proper-base-change (blueprint/chap-Proper-base-change/_section.md)
\section*{\texorpdfstring{4\quad Proper base change}{4 Proper base change}}
\phantomsection\label{bp1}%
\addcontentsline{toc}{section}{\texorpdfstring{4\quad Proper base change}{4 Proper base change}}
% text: prose-04-01
\phantomsection\label{bp2}%
The function–sheaf dictionary (Section \bpout{1.3}, Proposition \bpout{12.28}) turns the identity \(\#X_0(\mathbf F_q)=\sum_y\#f^{-1}(y)(\mathbf F_q)\) into a statement about sheaves, and for this the trace function of \(Rf_!K\) at a point \(y\) must be the sum of \(t_K\) over the fibre \(f^{-1}(y)\). This requires the stalks of \(Rf_!K\) to be the cohomology of the fibres, but nothing in the definition of the derived direct image says so.

% object: question-stalks-proper-direct-image (blueprint/chap-Proper-base-change/question-stalks-proper-direct-image/_section.md)
\phantomsection\label{bp3}%
\begin{bpstatement}{Question~4.1}{ (Stalks of higher direct images along a proper map)}
For \(f\) proper and \(F\) a torsion sheaf on \(X\), is the stalk \((R^qf_*F)_{\bar y}\) equal to \(H^q(X_{\bar y},F)\)?
\end{bpstatement}

% text: prose-04-02
\phantomsection\label{bp4}%
For a proper map \(f:X\to Y\) of locally compact Hausdorff spaces the answer is yes. Indeed, the sets \(f^{-1}(V)\), for \(V\) a neighbourhood of \(y\), form a basis of neighbourhoods of the compact fibre, and the cohomology of a compact subset is the colimit of the cohomology of its neighbourhoods.

In the étale topology, however, the stalk of \(R^qf_*F\) at a geometric point \(\bar y\) is the cohomology of \(X\times_YY_{(\bar y)}\), where \(Y_{(\bar y)}\) is the spectrum of the strict henselization of \(Y\) at \(\bar y\). This scheme is much larger than the fibre \(X_{\bar y}\). For instance, if \(f\) is the identity of a curve and \(\bar y\) lies over a closed point, then it is the spectrum of a strictly henselian local ring of dimension one, not a point.

Properness cannot be dropped. For example, let \(j:\mathbf G_m\hookrightarrow\mathbf A^1\) be the inclusion, \(i:\{0\}\hookrightarrow\mathbf A^1\) the origin, and \(A\) the constant sheaf on \(\mathbf G_m\) with value a nonzero finite abelian group. Then the sheaf \(i^*j_*A\) is nonzero, while the base change of \(\mathbf G_m\) to the origin is empty.

For proper \(f\), nevertheless, the answer is yes. More generally, pulling a proper family back to a new base and taking its cohomology can be done in either order. Precisely, for a cartesian square with \(f\) proper, \(g:Y'\to Y\) arbitrary and \(X'=X\times_YY'\), the base change map

\[
g^*Rf_*K\longrightarrow Rf'_*g'^*K
\]

is an isomorphism for every bounded below classical complex \(K\) with torsion cohomology sheaves (Theorem \hyperref[bp98]{4.33}). In particular, taking for \(g\) a geometric point gives the stalk formula (Corollary \hyperref[bp102]{4.34}). The theorem is due to Artin [\hyperref[bp123]{SGA4}, XII–XIII], and unlike smooth base change (Chapter \bpout{8}), it holds for torsion of every order, including torsion divisible by the residue characteristics.

The content of the theorem is already visible in degree zero, where it is a statement about connectedness. Indeed, sections of the constant sheaf \(\{0,1\}\) are open and closed subsets, and for a proper scheme over a henselian local ring, the open and closed subsets correspond to those of the closed fibre. When \(X\) is finite over a henselian local ring \(R\), this says that idempotents of the finite \(R\)-algebra lift from its reduction modulo the maximal ideal, which is one form of Hensel's lemma. Theorem \hyperref[bp24]{4.8} extends the statement to sections of every étale sheaf of sets. Higher degrees then reduce to degree zero and the projective line, and from \(\mathbf P^1\) the general case follows by formal steps, through \((\mathbf P^1)^n\), \(\mathbf P^n\) and Chow's lemma (Proposition \hyperref[bp94]{4.32}).

We make these calculations on the étale site, where the classical tools are available, and transport them to classical complexes on the pro-étale site by the comparison theorems of Chapter \bpout{2}. For \(\mathcal O_E\) we pass to the limit over the reductions modulo \(\pi_E^n\) (Proposition \hyperref[bp113]{4.38}). Since pullback need not preserve derived completeness, at this stage the statement for \(\mathcal O_E\) involves the completed inverse image. Once we know that \(Rp_*\) preserves constructibility (Proposition \hyperref[bp121]{4.40}), however, the completion is unnecessary, and we obtain the ordinary base change isomorphism on \(D^b_{\mathrm{ctf}}\) for every admissible coefficient ring.

We use proper base change three times later. First, it computes the stalks of \(Rf_!\) in the trace formula (Corollary \bpout{12.26}). Second, it makes \(Rf_!=Rp_*j_!\) independent of the compactification (Proposition \bpout{5.26}). Finally, it bounds the Tor amplitude of \(Rp_*\) with finite coefficients (Proposition \hyperref[bp106]{4.35}).

% text: situation-conventions
\phantomsection\label{bp5}%
\textbf{Convention.} For proper base change, schemes are arbitrary unless an additional hypothesis is stated. The theorem concerns the classical part of \(D^+(T_{\mathrm{pro}\acute{e}t},\Lambda)\), for a discrete commutative ring \(\Lambda\), but the intermediate calculations (acyclicity, curves, Chow's lemma) are made on \(T_{\acute{e}t}\), with abelian sheaves or modules over its discrete constant ring. There we write \(D^+(T_{\acute{e}t},\Lambda)\) for the bounded below derived category. Inverse image of étale sheaves along \(u:S\to T\) is written \(u^*:\Sh(T_{\acute{e}t})\to\Sh(S_{\acute{e}t})\), as for module sheaves; it is exact, and so is its extension to \(D^+(T_{\acute{e}t},\Lambda)\to D^+(S_{\acute{e}t},\Lambda)\). The derived direct image is written \(Ru_*\). Torsion always means torsion of the underlying abelian sheaves.

Proper base change holds for torsion coefficients of every characteristic. Accordingly, until Section \hyperref[bp104]{4.7}, \(\ell\) denotes an arbitrary prime, which may equal a residue characteristic. From that section on, we return to the standing finite-type hypotheses and the standing condition on \(\ell\).

% object: sec-The-base-change-map (blueprint/chap-Proper-base-change/sec-The-base-change-map/_section.md), content hidden
\subsection*{\texorpdfstring{4.1\quad The base change map}{4.1 The base change map}}
\phantomsection\label{bp6}%
\addcontentsline{toc}{subsection}{\texorpdfstring{4.1\quad The base change map}{4.1 The base change map}}
% text: prose-04-03
\phantomsection\label{bp7}%
Proper base change asserts that a specific map, from the pullback of the direct image to the direct image of the pullback, is an isomorphism. We define this map and describe it on sections, so that we can compute it on stalks.

From Definition \hyperref[bp38]{4.14} on, we fix the coefficient ring. Integral and rational coefficients are treated in Section \hyperref[bp108]{4.8}. Results for abelian sheaves apply to sheaves of modules, by the following lemma.

% object: lemma-cohomology-modules-underlying-abelian (blueprint/chap-Proper-base-change/sec-The-base-change-map/lemma-cohomology-modules-underlying-abelian/_section.md)
\phantomsection\label{bp8}%
\begin{bpstatementit}{Lemma~4.2}{ (Modules and abelian sheaves)}
For a sheaf of \(\Lambda\)-modules on a site, its cohomology and its higher direct images along a morphism of sites agree with those of its underlying abelian sheaf. The same holds for \(Ru_*\) of bounded below complexes.
\end{bpstatementit}

% proof: proof/lemma-cohomology-modules-underlying-abelian
\begin{proof}[{Proof}]\phantomsection\label{bp9}%
An injective sheaf of modules is acyclic for the section functor on every object of the site, although it need not be injective as an abelian sheaf [Stacks, \href{https://stacks.math.columbia.edu/tag/03FD}{03FD}]. Hence an injective resolution of modules computes the cohomology of the underlying abelian sheaf, on every object. Applying this to the objects over the target gives the statement for \(R^qu_*\), and bounded below resolutions give it for \(Ru_*\).
\end{proof}

% text: prose-04-04
\phantomsection\label{bp10}%
Consider a cartesian square

\[
\begin{CD}
X' @>{g'}>> X\\
@V{f'}VV @VV{f}V\\
Y' @>{g}>> Y.
\end{CD}
\tag{4.1}
\]

Since \(gf'=fg'\), the composition isomorphisms of inverse images yield an isomorphism \(f'^*g^*\simeq g'^*f^*\) on sheaves. Inverse images are exact, so the same isomorphism exists on bounded below derived categories.

% object: definition-underived-base-change-map (blueprint/chap-Proper-base-change/sec-The-base-change-map/definition-underived-base-change-map/_section.md)
\phantomsection\label{bp11}%
\begin{bpstatement}{Definition~4.3}{ (Underived base change map)}
Let \(F\) be a sheaf on \(X_{\acute{e}t}\). The \emph{underived base change map}

\[
g^*f_*F\longrightarrow f'_*g'^*F
\tag{4.2}
\]

is the adjoint, under \(f'^*\dashv f'_*\), of

\[
f'^*g^*f_*F
\simeq g'^*f^*f_*F
\xrightarrow{g'^*\varepsilon_f}g'^*F,
\]

where \(\varepsilon_f:f^*f_*\to\mathrm{id}\) is the counit.
\end{bpstatement}

% text: prose-04-05
\phantomsection\label{bp12}%
Equivalently, first apply the unit \(g^*f_*F\to f'_*f'^*g^*f_*F\), and then apply \(f'_*\) to the displayed composite. This construction is natural in \(F\), and it applies equally to sheaves of \(\Lambda\)-modules.

On sections the map has the following simple description. A section of \(g^*f_*F\) over an étale \(V\to Y'\) is represented by an étale \(W\to Y\), a map \(V\to
Y'\times_YW\), and a section over \(X\times_YW\). Given such data, the base change map restricts this section along

\[
X'\times_{Y'}V\longrightarrow X\times_YW,
\]

and it is the sheafification of this map of presheaves.

% object: definition-derived-base-change-map (blueprint/chap-Proper-base-change/sec-The-base-change-map/definition-derived-base-change-map/_section.md)
\phantomsection\label{bp13}%
\begin{bpstatement}{Definition~4.4}{ (Derived base change map)}
For \(K\in D^+(X_{\acute{e}t},\Lambda)\), the \emph{base change map}

\[
g^*Rf_*K\longrightarrow Rf'_*g'^*K
\tag{4.3}
\]

is the adjoint, under \(f'^*\dashv Rf'_*\), of

\[
f'^*g^*Rf_*K
\simeq g'^*f^*Rf_*K
\xrightarrow{g'^*\varepsilon_f}g'^*K,
\]

where \(\varepsilon_f:f^*Rf_*\to\mathrm{id}\) now denotes the derived counit. The morphism is natural in \(K\).
\end{bpstatement}

% text: prose-04-06
\phantomsection\label{bp14}%
The derived adjunction exists because inverse image is exact and its right adjoint preserves injectives. Note that the construction of the base change map uses only that the square commutes; the cartesian hypothesis is needed only when we ask whether the map is an isomorphism.

To compute (4.3), choose a bounded below injective resolution \(K\to I^\bullet\) and an injective replacement \(g'^*I^\bullet\to J^\bullet\). Then the base change map is represented by the morphism of complexes

\[
g^*f_*I^\bullet
\longrightarrow
f'_*g'^*I^\bullet\longrightarrow f'_*J^\bullet.
\tag{4.4}
\]

Here the first arrow is the underived base change map in each degree. By exactness of \(g^*\) the source is \(g^*Rf_*K\), and the last complex computes \(Rf'_*g'^*K\). The formula does not depend on the choices, since comparison maps between injective resolutions are unique up to homotopy. Moreover, it is the adjoint map of Definition \hyperref[bp13]{4.4}, because under adjunction the first arrow corresponds to the square isomorphism followed by the counit, and the replacement map is the canonical passage to the derived direct image.

Base change maps for composite squares are composites of base change maps.

% object: lemma-identity-pasting (blueprint/chap-Proper-base-change/sec-The-base-change-map/lemma-identity-pasting/_section.md)
\phantomsection\label{bp15}%
\begin{bpstatementit}{Lemma~4.5}{ (Identity and pasting)}
Let \(K\in D^+(X_{\acute{e}t},\Lambda)\). The base change map is the identity for an identity square, and is compatible with horizontal and vertical pasting of commutative squares, in particular of cartesian squares. More explicitly, if \(h:Y''\to Y'\) is a second base change, completed by maps \(h':X''\to X'\) and \(f'':X''\to Y''\) to a commutative square with \(f'h'=hf''\), the base change map for the pasted square, with bottom map \(gh\), is the composite

\[
h^*g^*Rf_*K
\longrightarrow
h^*Rf'_*g'^*K
\longrightarrow
Rf''_*h'^*g'^*K
\]

of \(h^*\) applied to the base change map of (4.1) and the base change map of the second square, under the canonical composition isomorphisms.

For composable maps \(X\xrightarrow{f}Y\xrightarrow{g}Z\) and a map \(h:Z'\to Z\), let \(g':Y'\to Z'\), \(h_Y:Y'\to Y\) and \(f':X'\to Y'\), \(h_X:X'\to X\) complete \(g\), \(h\) and \(f\), \(h_Y\) to commutative squares (for instance \(Y'=Y\times_ZZ'\) and \(X'=X\times_ZZ'\)). Then the base change map for the pasted square with left map \(g'f'\) and right map \(gf\) is the composite

\[
h^*Rg_*Rf_*K
\longrightarrow
Rg'_*h_Y^*Rf_*K
\longrightarrow
Rg'_*Rf'_*h_X^*K
\]

of the base change map for \(g\) at \(Rf_*K\) and \(Rg'_*\) applied to the base change map for \(f\) at \(K\).
\end{bpstatementit}

% proof: proof/lemma-identity-pasting
\begin{proof}[{Proof}]\phantomsection\label{bp16}%
Each base change map is a unit followed by a square isomorphism and a counit. When two squares are pasted, the square isomorphisms compose, and the triangle identities cancel the inner unit and counit. In the vertical case we also use that direct image preserves injectives, so that \(R(gf)_*\simeq Rg_*Rf_*\) is computed on the same resolution.
\end{proof}

% text: prose-04-07
\phantomsection\label{bp17}%
On the pro-étale site we use the same construction, but here \(\Lambda\) may be any coefficient ring pulled back from the point, not only a torsion ring. For such \(\Lambda\) inverse images are exact (Lemma \bpout{2.54}), and the adjunction \(f^*\dashv Rf_*\) exists on the unbounded derived categories.

% object: definition-pro-etale-base-change-map (blueprint/chap-Proper-base-change/sec-The-base-change-map/definition-pro-etale-base-change-map/_section.md)
\phantomsection\label{bp18}%
\begin{bpstatement}{Definition~4.6}{ (Pro-étale base change map)}
Let \(K\in D(X_{\mathrm{pro}\acute{e}t},\Lambda)\). The \emph{base change map}

\[
g^*Rf_*K\longrightarrow Rf'_*g'^*K
\tag{4.5}
\]

in \(D(Y'_{\mathrm{pro}\acute{e}t},\Lambda)\) is the adjoint transpose, for \(f'^*\dashv Rf'_*\), of the composite

\[
f'^*g^*Rf_*K\simeq g'^*f^*Rf_*K
\xrightarrow{g'^*\varepsilon_f}g'^*K,
\]

where \(\varepsilon_f:f^*Rf_*\to\mathrm{id}\) is the counit of \(f^*\dashv Rf_*\).
\end{bpstatement}

% text: prose-04-08
\phantomsection\label{bp19}%
Equivalently, with \(\eta_{f'}\) the unit of \(f'^*\dashv Rf'_*\), it is the composite

\[
g^*Rf_*K\xrightarrow{\eta_{f'}}
Rf'_*f'^*g^*Rf_*K\simeq Rf'_*g'^*f^*Rf_*K
\xrightarrow{Rf'_*g'^*\varepsilon_f}Rf'_*g'^*K.
\]

This is the pro-étale form of the base change map of Definition \hyperref[bp13]{4.4}. It is defined for every cartesian square (4.1), and the same construction applies to any commutative square. Since proper and smooth base change are proved on the étale site and then transported to the pro-étale site, we need to know that the two maps correspond.

% object: lemma-naturality-pasting-comparison (blueprint/chap-Proper-base-change/sec-The-base-change-map/lemma-naturality-pasting-comparison/_section.md)
\phantomsection\label{bp20}%
\begin{bpstatementit}{Lemma~4.7}{ (Naturality, pasting and comparison)}
The maps (4.5) form a natural transformation \(g^*Rf_*\to Rf'_*g'^*\). It is the identity for an identity square, and for a horizontal or vertical pasting of two commutative squares the base change map of the pasted square is the composite of the two base change maps, in the form displayed in Lemma \hyperref[bp15]{4.5}. For a discrete ring \(\Lambda\) and \(K\in D^+(X_{\acute{e}t},\Lambda)\), with \(f\) and \(f'\) qcqs, the comparison isomorphisms of Theorem \bpout{2.38} identify the base change map of \(\nu_X^*K\) with the inverse image on \(Y'_{\mathrm{pro}\acute{e}t}\) of the étale base change map of \(K\).
\end{bpstatementit}

% proof: proof/lemma-naturality-pasting-comparison
\begin{proof}[{Proof}]\phantomsection\label{bp21}%
Every assertion holds because the base change map is assembled from units, counits and composition isomorphisms for inverse images, and each of the operations in the lemma respects these ingredients. Naturality in \(K\) is immediate, since the counit \(\varepsilon_f\) is a natural transformation.

By definition, the base change map is the adjoint transpose of a counit preceded by the composition isomorphism for inverse images. Since the counit for a composite adjunction is the composite of the two counits, taking adjoint transposes proves pasting, in either order of composing the two directions of a rectangle. Similarly, the assertion about identity squares follows from the triangle identities, as in Lemma \hyperref[bp15]{4.5}. Finally, the comparisons of sites of Theorem \bpout{2.38} identify the composition maps for inverse images and the units and counits (Lemma \bpout{2.39}), and hence also the base change maps constructed from them.
\end{proof}


% object: sec-Base-change-in-degree-zero (blueprint/chap-Proper-base-change/sec-Base-change-in-degree-zero/_section.md), content hidden
\subsection*{\texorpdfstring{4.2\quad Base change in degree zero}{4.2 Base change in degree zero}}
\phantomsection\label{bp22}%
\addcontentsline{toc}{subsection}{\texorpdfstring{4.2\quad Base change in degree zero}{4.2 Base change in degree zero}}
% text: prose-04-09
\phantomsection\label{bp23}%
For a proper morphism, the underived base change map is an isomorphism for every étale sheaf of sets, with no torsion hypothesis. The reason is that, for a proper scheme over a henselian pair, restriction of sections to the closed fibre is bijective (Theorem \hyperref[bp24]{4.8}).

Recall that a pair \((R,I)\), with \(I\) an ideal of the ring \(R\), is \emph{henselian} if \(I\) is contained in the Jacobson radical of \(R\) and Hensel's lemma holds for it: every factorization modulo \(I\) of a monic polynomial into coprime monic factors lifts to a factorization in \(R\) [Stacks, \href{https://stacks.math.columbia.edu/tag/09XE}{09XE}]. A local ring with its maximal ideal is a henselian pair exactly when it is henselian. For example, a complete local ring is henselian.

% object: theorem-comparison-over-henselian-pair-degree (blueprint/chap-Proper-base-change/sec-Base-change-in-degree-zero/theorem-comparison-over-henselian-pair-degree/_section.md)
\phantomsection\label{bp24}%
\begin{bpstatementit}{Theorem~4.8}{ (Henselian comparison in degree zero)}
Let \((R,J)\) be a henselian pair, let \(Z\to\operatorname{Spec}R\) be proper, and let \(i:Z_0=Z\times_R\operatorname{Spec}(R/J)\hookrightarrow Z\). For every sheaf of sets \(F\) on \(Z_{\acute{e}t}\), restriction is a bijection

\[
\Gamma(Z,F)\xrightarrow{\sim}\Gamma(Z_0,i^*F).
\]

The comparison is natural in \(F\).
\end{bpstatementit}

% text: prose-04-10
\phantomsection\label{bp25}%
The theorem is [Stacks, \href{https://stacks.math.columbia.edu/tag/0A0C}{0A0C}]. At heart it is a statement about connectedness: every irreducible closed subset of a scheme proper over a henselian pair meets the closed fibre in a nonempty connected subset. This follows from Stein factorization and the connectedness properties of henselian pairs [Stacks, \href{https://stacks.math.columbia.edu/tag/0A0B}{0A0B}]; over a general affine base, Stein factorization only produces an algebra integral over \(R\), but that suffices. A topological argument on spectral spaces then turns connectedness into bijectivity of restriction for Zariski sheaves.

Passing to étale sheaves uses finite covers, which are harmless here: a finite cover of \(Z\) remains proper over \(R\), so the connectedness statement also holds after every finite cover. Given an étale covering of \(Z\), we choose a finite surjection \(Z'\to Z\) that factors through the covering Zariski locally on \(Z'\), and sections then descend along \(Z'\times_ZZ'\rightrightarrows Z'\) [Stacks, \href{https://stacks.math.columbia.edu/tag/09ZF}{09ZF}].

% object: example-henselian-needed (blueprint/chap-Proper-base-change/sec-Base-change-in-degree-zero/example-henselian-needed/_section.md)
\phantomsection\label{bp26}%
\begin{bpstatement}{Example~4.9}{ (Why the pair must be henselian)}
Let \(k\) be a field of characteristic \(\neq2\) and let \(R=k[t]_{(t)}\), the local ring of \(\mathbf A^1_k\) at the origin, with maximal ideal \(J=(t)\). The finite \(R\)-algebra \(B=R[x]/(x^2-1-t)\) is a domain, because \(1+t\) is not a square in \(k(t)\). Hence \(Z=\operatorname{Spec}B\) is connected and \(\Gamma(Z,\mathbf Z/2)=\mathbf Z/2\). The closed fibre, on the other hand, is \(\operatorname{Spec}k[x]/(x^2-1)\), which consists of the two points \(x=\pm1\), and so \(\Gamma(Z_0,\mathbf Z/2)=(\mathbf Z/2)^2\). Thus the conclusion of the theorem fails for \((R,J)\). The reason is that \((R,J)\) is not henselian: the factorization \(x^2-1=(x-1)(x+1)\) modulo \(t\) does not lift. In the henselization, by contrast, and in \(k[[t]]\), the element \(1+t\) has the square root \(1+t/2-t^2/8+\cdots\). There the factorization lifts, and \(Z\) splits into two copies of the base, so both sides become \((\mathbf Z/2)^2\).
\end{bpstatement}

% text: prose-04-11
\phantomsection\label{bp27}%
The stalks of \(f_*F\) are sections over geometric fibres. Indeed, a stalk is the set of sections over the strict localization, and the henselian comparison identifies these with the sections over the closed fibre.

% object: lemma-sections-geometric-fibre (blueprint/chap-Proper-base-change/sec-Base-change-in-degree-zero/lemma-sections-geometric-fibre/_section.md)
\phantomsection\label{bp28}%
\begin{bpstatementit}{Lemma~4.10}{ (Sections over a geometric fibre)}
Let \(f:X\to Y\) be proper, let \(F\) be an étale sheaf of sets on \(X\), and let \(\bar y:\operatorname{Spec}\Omega\to Y\) be a geometric point. Then restriction is a bijection

\[
(f_*F)_{\bar y}\xrightarrow{\sim}
\Gamma(X_{\bar y},F|_{X_{\bar y}}).
\tag{4.6}
\]
\end{bpstatementit}

% proof: proof/lemma-sections-geometric-fibre
\begin{proof}[{Proof}]\phantomsection\label{bp29}%
Let \(y\in Y\) be the image of \(\bar y\), let \(Y_{(\bar y)}=\operatorname{Spec}\mathcal O^{\mathrm{sh}}_{Y,\bar y}\) and \(X_{(\bar y)}=X\times_YY_{(\bar y)}\). In degree zero the stalk is computed on the strict localization, that is,

\[
(f_*F)_{\bar y}\simeq\Gamma(X_{(\bar y)},F|_{X_{(\bar y)}}).
\]

For quasi-compact quasi-separated \(f\) this holds by continuity of étale sections ([Stacks, \href{https://stacks.math.columbia.edu/tag/03Q9}{03Q9}], whose proof in degree zero applies to sheaves of sets). To recall the argument: by definition, the stalk \((f_*F)_{\bar y}\) is the filtered colimit of \(\Gamma(X\times_YV,F|)\) over pointed affine étale neighbourhoods \(V\) of \(\bar y\), where \(F|\) denotes the inverse image of \(F\). Since the limit of these neighbourhoods is \(Y_{(\bar y)}\), there is a natural map from the colimit to the sections over \(X_{(\bar y)}\), and we must see that it is bijective. A section over the limit is represented by finitely many covering objects and finitely many agreement conditions, because \(X_{(\bar y)}\) is quasi-compact and quasi-separated. Since all these data are of finite presentation, they descend to one neighbourhood. This proves surjectivity, and the same descent for two representatives proves injectivity.

The local ring of \(Y_{(\bar y)}\) is strictly henselian, with separably closed residue field \(\kappa(y)^{\mathrm{sep}}\subseteq\Omega\). Thus \(\mathcal O^{\mathrm{sh}}_{Y,\bar y}\) with its maximal ideal is a henselian pair, and \(X_{(\bar y)}\) is proper over it with closed fibre \(X_{\kappa(y)^{\mathrm{sep}}}\). Theorem \hyperref[bp24]{4.8} therefore identifies the sections on the right with those over \(X_{\kappa(y)^{\mathrm{sep}}}\). These in turn equal \(\Gamma(X_{\bar y},F|_{X_{\bar y}})\), because extending a separably closed ground field from \(\kappa(y)^{\mathrm{sep}}\) to \(\Omega\) does not change the sections of the pullback of an étale sheaf [Stacks, \href{https://stacks.math.columbia.edu/tag/0A3I}{0A3I}].
\end{proof}

% object: proposition-proper-base-change-degree-zero (blueprint/chap-Proper-base-change/sec-Base-change-in-degree-zero/proposition-proper-base-change-degree-zero/_section.md)
\phantomsection\label{bp30}%
\begin{bpstatementit}{Proposition~4.11}{ (Proper base change in degree zero)}
For (4.1) with \(f\) proper, and for every étale sheaf of sets \(F\) on \(X\), the underived base change map (4.2)

\[
g^*f_*F\longrightarrow f'_*g'^*F
\]

is an isomorphism.
\end{bpstatementit}

% proof: proof/proposition-proper-base-change-degree-zero
\begin{proof}[{Proof}]\phantomsection\label{bp31}%
We show that both stalks are the sections over the same geometric fibre. To this end, choose \(\bar y':\operatorname{Spec}\Omega\to Y'\) and use the composite \(\bar y=g\bar y'\) as the geometric point of \(Y\). Since properness is preserved by base change, Lemma \hyperref[bp28]{4.10} applies to \(f'\) as well as to \(f\). Moreover, the cartesian square identifies

\[
X'_{\bar y'}=X\times_Y\operatorname{Spec}\Omega=X_{\bar y}
\]

and identifies the two inverse image sheaves on it. Thus the stalks of both sides of the desired map are the same set of sections. It remains to see that the map is the identity. The underived base change map sends a section represented over \(X\times_YV\) to its restriction to the base change of the neighbourhood, and hence to its restriction to this same fibre. In other words, the map on stalks is the identity under the displayed identifications. Since geometric stalks detect isomorphisms, this proves the assertion.
\end{proof}


% object: sec-Torsion-sheaves-and-acyclicity (blueprint/chap-Proper-base-change/sec-Torsion-sheaves-and-acyclicity/_section.md), content hidden
\subsection*{\texorpdfstring{4.3\quad Torsion sheaves and acyclicity}{4.3 Torsion sheaves and acyclicity}}
\phantomsection\label{bp32}%
\addcontentsline{toc}{subsection}{\texorpdfstring{4.3\quad Torsion sheaves and acyclicity}{4.3 Torsion sheaves and acyclicity}}
% text: prose-04-12
\phantomsection\label{bp33}%
In degree zero we needed no hypothesis on the sheaf. In higher degrees we do: proper base change fails for the constant sheaf \(\mathbf Z\).

% object: example-pbc-fails-integral (blueprint/chap-Proper-base-change/sec-Torsion-sheaves-and-acyclicity/example-pbc-fails-integral/_section.md)
\phantomsection\label{bp34}%
\begin{bpstatement}{Example~4.12}{ (Proper base change fails for $\mathbf Z$)}
Let \(R\) be a strictly henselian discrete valuation ring with uniformizer \(t\) and algebraically closed residue field of characteristic \(\neq2\), and let \(X\subset\mathbf P^2_R\) be the closure of the affine curve \(y^2=x^3-x^2+t\). Then \(X\) is regular (at the origin of the closed fibre the equation has the linear term \(t\)), while its closed fibre \(X_s\) is the nodal cubic \(y^2=x^2(x-1)\). Let \(f:X\to\operatorname{Spec}R\) be the structure map. Since \(R\) is strictly henselian, the stalk of \(R^1f_*\mathbf Z\) at the closed point is \(H^1(X,\mathbf Z)\).

This group is zero. Indeed, \(X\) is normal and connected, so \(\mathbf Z\)-torsors on \(X\) correspond to continuous homomorphisms \(\pi_1(X)\to\mathbf Z\), and the image of such a homomorphism is a compact subgroup of \(\mathbf Z\), hence zero. On the closed fibre, however, \(H^1(X_s,\mathbf Z)=\mathbf Z\), generated by the loop through the node. This means that the base change map from the stalk of \(R^1f_*\mathbf Z\) at the closed point to \(H^1(X_s,\mathbf Z)\) is not surjective.

With \(\mathbf Z/n\) in place of \(\mathbf Z\), the class of the loop is realized by the cyclic étale cover of \(X_s\) of degree \(n\) obtained by unwinding the node, and because \(R\) is henselian, this cover lifts to \(X\). This is the phenomenon the torsion hypothesis captures.
\end{bpstatement}

% text: prose-04-13
\phantomsection\label{bp35}%
We therefore restrict to torsion sheaves.

% object: definition-torsion-sheaves (blueprint/chap-Proper-base-change/sec-Torsion-sheaves-and-acyclicity/definition-torsion-sheaves/_section.md)
\phantomsection\label{bp36}%
\begin{bpstatement}{Definition~4.13}{ (Torsion sheaves)}
A sheaf \(F\) of abelian groups, or of \(\Lambda\)-modules, is \emph{torsion} if every geometric stalk is a torsion abelian group. Equivalently,

\[
F=\varinjlim_{n\geq1}F[n],\qquad F[n]=\ker(n:F\to F),
\]

where the positive integers are ordered by divisibility.
\end{bpstatement}

% text: prose-04-14
\phantomsection\label{bp37}%
If \(\Lambda\) is torsion as an abelian group, it is killed by the additive order of \(1\), so every sheaf of \(\Lambda\)-modules is torsion. The two descriptions of torsion in Definition \hyperref[bp36]{4.13} agree. Indeed, the colimit is taken in sheaves, and geometric stalks commute with colimits and finite limits and detect isomorphisms, so the equivalence reduces to the case of abelian groups.

Torsion is a local condition on sections: every section is killed by an integer after passing to an étale covering. However, the integer may vary between members of the covering. For example, the constant sheaf \(\mathbf Q/\mathbf Z\) on a nonempty scheme is torsion, but on the disjoint union of countably many copies of the scheme its sections include \((1/2^n)_{n\geq1}\), which has infinite order.

% object: definition-torsion-pbc-property (blueprint/chap-Proper-base-change/sec-Torsion-sheaves-and-acyclicity/definition-torsion-pbc-property/_section.md)
\phantomsection\label{bp38}%
\begin{bpstatement}{Definition~4.14}{ (The PBC property)}
For a morphism \(f:X\to Y\), write \(\mathrm{PBC}_{\Lambda}(f)\) if the base change map (4.3) is an isomorphism for every base change \(g:Y'\to Y\) and every \(K\in D^+(X_{\acute{e}t},\Lambda)\) with torsion cohomology sheaves.
\end{bpstatement}

% text: prose-04-15
\phantomsection\label{bp39}%
From now until the end of Section \hyperref[bp78]{4.5}, \(\Lambda\) denotes a commutative ring that is torsion as an abelian group, unless another ring is specified. For such \(\Lambda\) and proper \(f\) we prove \(\mathrm{PBC}_\Lambda(f)\) in Propositions \hyperref[bp76]{4.26} and \hyperref[bp94]{Proposition 4.32}, and Theorem \hyperref[bp98]{4.33} then removes the condition on \(\Lambda\).

To pass from degree zero to higher degrees we apply the degree-zero result to an injective resolution \(F\to I^\bullet\) of a sheaf \(F\) on \(X\). By Proposition \hyperref[bp30]{4.11} the base change map is an isomorphism \(g^*f_*I^m\simeq f'_*g'^*I^m\) in each degree. Here the complex \(g^*f_*I^\bullet\) computes \(g^*Rf_*F\). The complex \(f'_*g'^*I^\bullet\), on the other hand, is \(f'_*\) applied to a resolution of \(g'^*F\) that need not be injective, since \(g'^*\) does not preserve injectives. It still computes \(Rf'_*g'^*F\) if every \(g'^*I^m\) is \(f'_*\)-acyclic. We call this condition injective acyclicity, and Theorem \hyperref[bp45]{4.17} shows that it is equivalent to proper base change. Its advantage is that, unlike the base change statement, injective acyclicity is a vanishing condition on single sheaves.

% object: definition-injective-acyclicity (blueprint/chap-Proper-base-change/sec-Torsion-sheaves-and-acyclicity/definition-injective-acyclicity/_section.md)
\phantomsection\label{bp40}%
\begin{bpstatement}{Definition~4.15}{ (Injective acyclicity)}
For a ring \(\Lambda\), write \(\mathrm{IA}_{\Lambda}(f)\) if for every cartesian square (4.1) and every injective sheaf \(I\) of \(\Lambda\)-modules on \(X_{\acute{e}t}\) which is torsion, one has

\[
R^qf'_*g'^*I=0\qquad(q>0).
\]
\end{bpstatement}

% text: prose-04-16
\phantomsection\label{bp41}%
In general \(g'^*I\) is not injective, and injective acyclicity only asks that it be \(f'_*\)-acyclic. By the next lemma this suffices to compute \(Rf'_*\) of the pullback of an injective resolution term by term.

% object: lemma-acyclic-resolutions (blueprint/chap-Proper-base-change/sec-Torsion-sheaves-and-acyclicity/lemma-acyclic-resolutions/_section.md)
\phantomsection\label{bp42}%
\begin{bpstatementit}{Lemma~4.16}{ (Acyclic resolutions)}
Let \(v:\mathcal A\to\mathcal B\) be a left exact additive functor between abelian categories with enough injectives. If \(F\to M^\bullet\) is a resolution and each \(M^m\) is \(v\)-acyclic, then

\[
\mathcal H^q(vM^\bullet)\simeq R^qvF.
\]

More generally, for a bounded below complex \(M^\bullet\) of \(v\)-acyclic objects, the canonical comparison \(vM^\bullet\to Rv(M^\bullet)\) is an isomorphism in \(D^+(\mathcal B)\).
\end{bpstatementit}

% proof: proof/lemma-acyclic-resolutions
\begin{proof}[{Proof}]\phantomsection\label{bp43}%
This is Leray's acyclicity lemma [Stacks, \href{https://stacks.math.columbia.edu/tag/015E}{015E}], and the first assertion is the case of a resolution of \(F\). We recall the argument, which compares \(M^\bullet\) with an injective replacement through the cone.

In a bounded below exact complex of \(v\)-acyclic objects the cycles are again \(v\)-acyclic, by induction along the short exact sequences of cycles, so the complex stays exact after applying \(v\). We apply this to the cone of a quasi-isomorphism \(M^\bullet\to J^\bullet\) into a bounded below complex of injectives. This cone is bounded below and exact, and its terms are \(v\)-acyclic, being finite direct sums of terms of \(M^\bullet\) and of \(J^\bullet\). It follows that \(vM^\bullet\to vJ^\bullet\) is a quasi-isomorphism, and the target computes \(Rv(M^\bullet)\).
\end{proof}

% text: prose-04-17
\phantomsection\label{bp44}%
To pass from sheaves killed by an integer to all torsion sheaves, we use continuity of étale cohomology: if \(u:T\to S\) is quasi-compact and quasi-separated, then

\[
\varinjlim_i R^qu_*F_i\xrightarrow{\sim}
R^qu_*\bigl(\varinjlim_i F_i\bigr)\qquad(q\geq0)
\tag{4.7}
\]

for every filtered system of abelian étale sheaves on \(T\). The same holds after any base change of \(u\) [Stacks, \href{https://stacks.math.columbia.edu/tag/09Z1}{09Z1}]. Indeed, \(R^qu_*F\) is the sheafification of \(V\mapsto H^q(T\times_SV,F|_{T\times_SV})\) on affine étale \(V\to S\), and the étale cohomology of the quasi-compact and quasi-separated scheme \(T\times_SV\) commutes with filtered colimits of abelian sheaves [Stacks, \href{https://stacks.math.columbia.edu/tag/03Q5}{03Q5}].

The following criterion reduces proper base change in higher degrees to injective acyclicity, a vanishing statement that can be checked on curves.

% object: theorem-pbc-injective-acyclicity (blueprint/chap-Proper-base-change/sec-Torsion-sheaves-and-acyclicity/theorem-pbc-injective-acyclicity/_section.md)
\phantomsection\label{bp45}%
\begin{bpstatementit}{Theorem~4.17}{ (PBC and injective acyclicity)}
Let \(f:X\to Y\) be proper. The following are equivalent:

\begin{itemize}
\item (i) \(\mathrm{PBC}_\Lambda(f)\) holds.
\item (ii) For every sheaf \(F\) of \(\Lambda\)-modules and every base change (4.1), the maps induced by (4.3) on cohomology sheaves \[
  g^*R^qf_*F\longrightarrow R^qf'_*g'^*F
  \] are isomorphisms for all \(q\geq0\).
\item (iii) \(\mathrm{IA}_\Lambda(f)\) holds.
\end{itemize}
\end{bpstatementit}

% proof: proof/theorem-pbc-injective-acyclicity
\begin{proof}[{Proof}]\phantomsection\label{bp46}%
We prove (i)\(\Rightarrow\)(ii)\(\Rightarrow\)(iii)\(\Rightarrow\)(ii)\(\Rightarrow\)(i). The step (iii)\(\Rightarrow\)(ii) is the main one: it applies the degree-zero isomorphism termwise to an injective resolution.

\emph{(i)\(\Rightarrow\)(ii)\(\Rightarrow\)(iii).} Condition (ii) is the case of (i) where \(F\) sits in degree zero. Suppose (ii) holds, and let \(I\) be an injective sheaf of \(\Lambda\)-modules. Then \(R^qf'_*g'^*I\simeq g^*R^qf_*I\) by (ii). The right hand side vanishes for \(q>0\) because \(I\) is injective, and this is (iii).

\emph{(iii)\(\Rightarrow\)(ii).} Let \(F\) be a sheaf of \(\Lambda\)-modules. We compare the two sides of the base change map through an injective resolution \(F\to I^\bullet\) by sheaves of \(\Lambda\)-modules. The complex \(f'_*g'^*I^\bullet\) computes \(Rf'_*g'^*F\). Indeed, since \(g'^*\) is exact, \(g'^*I^\bullet\) is a resolution of \(g'^*F\), and its terms are \(f'_*\)-acyclic by (iii), so Lemma \hyperref[bp42]{4.16} applies. Similarly, the complex \(g^*f_*I^\bullet\) computes \(g^*Rf_*F\), because inverse image is exact. Between the two, Proposition \hyperref[bp30]{4.11} gives an isomorphism

\[
g^*f_*I^m\xrightarrow{\sim} f'_*g'^*I^m
\]

in each degree, and these form an isomorphism of complexes by naturality. By (4.4) this isomorphism is the derived base change map for \(F\), which proves (ii).

\emph{(ii)\(\Rightarrow\)(i).} Let \(K\in D^+(X_{\acute{e}t},\Lambda)\). For bounded \(K\) we argue by induction on the cohomological length. A complex of length one is a shifted sheaf, which (ii) covers, while a longer one is, by a canonical truncation triangle, an extension of two complexes of smaller length. Since the base change map is a morphism of triangles, the five lemma on cohomology sheaves shows that it is an isomorphism for \(K\). An unbounded \(K\) reduces to the bounded case, because a given cohomology sheaf only depends on a bounded truncation. More precisely, fix a degree \(m\) and choose \(N\geq m\), so that \(\tau_{\leq N}K\) is bounded and has the desired comparison. Both functors \(g^*Rf_*\) and \(Rf'_*g'^*\) send \(D^{\geq N+1}(X_{\acute{e}t},\Lambda)\) into \(D^{\geq N+1}(Y'_{\acute{e}t},\Lambda)\), since inverse images are exact and derived direct images are left \(t\)-exact. Applying them to

\[
\tau_{\leq N}K\longrightarrow K\longrightarrow
\tau_{\geq N+1}K\xrightarrow{+1}
\]

we see that the comparison for \(K\) in degree \(m\) is the comparison for \(\tau_{\leq N}K\), hence an isomorphism. As \(m\) was arbitrary, (i) follows.
\end{proof}


% object: sec-The-projective-line (blueprint/chap-Proper-base-change/sec-The-projective-line/_section.md), content hidden
\subsection*{\texorpdfstring{4.4\quad The projective line}{4.4 The projective line}}
\phantomsection\label{bp47}%
\addcontentsline{toc}{subsection}{\texorpdfstring{4.4\quad The projective line}{4.4 The projective line}}
% text: prose-04-18
\phantomsection\label{bp48}%
Let \(\pi:\mathbf P^1_S\to S\) be the projection. A geometric stalk of \(R^q\pi_*F\) is the cohomology of \(\mathbf P^1\) over a strictly henselian local ring, and proper base change for \(\pi\) says that it is the cohomology of the geometric fibre, and that this group does not change under base change.

% object: subsec-Henselian-pairs-and-curves (blueprint/chap-Proper-base-change/sec-The-projective-line/subsec-Henselian-pairs-and-curves/_section.md), content hidden
\subsubsection*{\texorpdfstring{4.4.1\quad Henselian pairs and curves}{4.4.1 Henselian pairs and curves}}
\phantomsection\label{bp49}%
\addcontentsline{toc}{subsubsection}{\texorpdfstring{4.4.1\quad Henselian pairs and curves}{4.4.1 Henselian pairs and curves}}
% text: prose-04-19
\phantomsection\label{bp50}%
The model case is a complete discrete valuation ring \(R\). Think of \(\operatorname{Spec}R\) as a small disc around its closed point, and of \(\mathbf P^1_R\) as a family of projective lines over the disc. In topology such a family retracts onto the central fibre, and we expect that, for torsion sheaves, the cohomology of \(\mathbf P^1_R\) is that of its closed fibre. Besides the comparison in degree zero of Theorem \hyperref[bp24]{4.8}, the proof uses two theorems about henselian pairs from the Stacks Project. The first is the affine analogue of what we want to prove.

% object: theorem-gabber-affine-proper-base-change (blueprint/chap-Proper-base-change/sec-The-projective-line/subsec-Henselian-pairs-and-curves/theorem-gabber-affine-proper-base-change/_section.md)
\phantomsection\label{bp51}%
\begin{bpstatementit}{Theorem~4.18}{ (Gabber's affine analogue of proper base change)}
Let \((B,J)\) be a henselian pair and let \(G\) be a torsion abelian sheaf on \(\operatorname{Spec}(B)_{\acute{e}t}\). Then restriction induces isomorphisms

\[
H^q(\operatorname{Spec}B,G)
\xrightarrow{\sim}
H^q(\operatorname{Spec}(B/J),G|_{\operatorname{Spec}(B/J)})
\qquad(q\geq0).
\]
\end{bpstatementit}

% text: prose-04-20
\phantomsection\label{bp52}%
This is [Stacks, \href{https://stacks.math.columbia.edu/tag/09ZI}{09ZI}]; Gabber proves it by an effacement argument on affines. The second fact concerns finite étale covers of a proper scheme. In degree one it says that \(\mathbf Z/n\)-torsors, like sections, are determined by the closed fibre.

% object: theorem-finite-etale-covers-proper-henselian (blueprint/chap-Proper-base-change/sec-The-projective-line/subsec-Henselian-pairs-and-curves/theorem-finite-etale-covers-proper-henselian/_section.md)
\phantomsection\label{bp53}%
\begin{bpstatementit}{Theorem~4.19}{ (Finite étale covers over a proper henselian pair)}
Let \((B,J)\) be a henselian pair, let \(Y\) be proper over \(B\), and put \(Y_0=Y\times_{\operatorname{Spec}B}\operatorname{Spec}(B/J)\). Then restriction is an equivalence between the categories of finite étale schemes over \(Y\) and over \(Y_0\),

\[
\textit{F\'Et}_Y\xrightarrow{\sim}\textit{F\'Et}_{Y_0}.
\]

Consequently, for a finite abelian group \(M\),

\[
H^1(Y,\underline M)\xrightarrow{\sim}H^1(Y_0,\underline M).
\tag{4.8}
\]
\end{bpstatementit}

% proof: proof/theorem-finite-etale-covers-proper-henselian
\begin{proof}[{Proof}]\phantomsection\label{bp54}%
The equivalence is [Stacks, \href{https://stacks.math.columbia.edu/tag/0GS2}{0GS2}], proved by algebraization, approximation, and descent of finite presentation. For the consequence, note that the groups in (4.8) classify \(M\)-torsors, which are finite étale schemes with an action of \(M\) satisfying the equations of a simply transitive action. The equivalence transports these schemes, their actions and the equations, so it induces a bijection on isomorphism classes of \(M\)-torsors.
\end{proof}

% text: prose-04-21
\phantomsection\label{bp55}%
The last step of the proof, the change of separably closed ground field, uses the following facts about curves over an algebraically closed field, which we also need later. For torsion divisible by the characteristic, proper and affine curves behave differently, which is why some of the assertions require \(C\) to be proper.

% object: theorem-cohomology-curves (blueprint/chap-Proper-base-change/sec-The-projective-line/subsec-Henselian-pairs-and-curves/theorem-cohomology-curves/_section.md)
\phantomsection\label{bp56}%
\begin{bpstatementit}{Theorem~4.20}{ (Cohomology of curves)}
Let \(k\) be algebraically closed, let \(C\) be a separated scheme of finite type over \(k\) with \(\dim C\leq1\), and let \(F\) be a classical torsion abelian sheaf on \(C_{\mathrm{pro}\acute{e}t}\). With cohomology taken on \(C_{\mathrm{pro}\acute{e}t}\), the following assertions hold:

\begin{itemize}
\item \((C_{1})\) \(H^q(C,F)=0\) for \(q>2\).
\item \((C_{2})\) If \(C\) is affine, \(H^q(C,F)=0\) for \(q>1\).
\item \((C_{3})\) If \(\operatorname{char}k=p>0\) and \(F\) is \(p\)-primary, \(H^q(C,F)=0\) for \(q>1\).
\item \((C_{4})\) If \(F\) is constructible with finite stalks and prime to \(\operatorname{char}k\), all \(H^q(C,F)\) are finite.
\item \((C_{5})\) If \(C\) is proper and \(F\) is constructible with finite stalks, all \(H^q(C,F)\) are finite.
\item \((C_{6})\) If \(k'/k\) is an extension of algebraically closed fields and \(F\) is prime to \(\operatorname{char}k\), the pullback maps \(H^q(C,F)\to H^q(C_{k'},F|)\) are isomorphisms.
\item \((C_{7})\) If \(C\) is proper, the maps in \((C_6)\) are isomorphisms for every torsion \(F\).
\item \((C_{8})\) For any open \(U\subseteq C\), restriction \(H^2(C,F)\to H^2(U,F|_U)\) is surjective.
\end{itemize}
\end{bpstatementit}

% text: prose-04-22
\phantomsection\label{bp57}%
In characteristic zero, the condition that \(F\) be prime to the characteristic in \((C_4)\) and \((C_6)\) is automatic.

% proof: proof/theorem-cohomology-curves
\begin{proof}[{Proof of Theorem 4.20}]\phantomsection\label{bp58}%
The proof is in Appendix \bpout{D}. For a smooth proper connected curve and \(n\) invertible in \(k\), one computes \(H^q(C,\mu_n)\) with the Kummer sequence and the Jacobian (Appendix \bpout{C}), while for \(p\)-primary coefficients one uses the Artin–Schreier sequence instead. The case of general torsion sheaves is then reduced to these cases by normalization, removal of finitely many points and a finite étale cover of degree prime to the torsion order.
\end{proof}

% text: prose-04-23
\phantomsection\label{bp59}%
The final step for the projective line compares cohomology over the residue fields of two strict localizations, which are separably closed but in general not algebraically closed. For this we need the following form of \((C_7)\).

% object: lemma-extension-separably-closed-ground-field (blueprint/chap-Proper-base-change/sec-The-projective-line/subsec-Henselian-pairs-and-curves/lemma-extension-separably-closed-ground-field/_section.md)
\phantomsection\label{bp60}%
\begin{bpstatementit}{Lemma~4.21}{ (Separably closed field extension)}
Let \(k\subseteq k'\) be separably closed fields, let \(C\) be proper over \(k\) with \(\dim C\leq1\), and let \(F\) be a classical torsion abelian sheaf on \(C_{\mathrm{pro}\acute{e}t}\). Then the natural map

\[
H^q(C,F)\xrightarrow{\sim}H^q(C_{k'},F|)
\qquad(q\geq0)
\]

is an isomorphism.
\end{bpstatementit}

% proof: proof/lemma-extension-separably-closed-ground-field
\begin{proof}[{Proof}]\phantomsection\label{bp61}%
A purely inseparable extension does not change the étale topos, which allows us to replace \(k\) and \(k'\) by their algebraic closures and then apply \((C_7)\).

By the classical comparison we may compute on the étale site. Choose compatible algebraic closures \(\bar k\subseteq\bar k'\). Since \(\bar k/k\) and \(\bar k'/k'\) are purely inseparable, \(C_{\bar k}\to C\) and \(C_{\bar k'}\to C_{k'}\) are universal homeomorphisms, and hence they induce equivalences of small étale topoi [Stacks, \href{https://stacks.math.columbia.edu/tag/03SI}{03SI}]. This reduces the claim to the pullback from \(C_{\bar k}\) to \(C_{\bar k'}\), which is an isomorphism by \((C_7)\) of Theorem \hyperref[bp56]{4.20}. Since all these maps are pullback maps, the comparison square commutes [Stacks, \href{https://stacks.math.columbia.edu/tag/0A5E}{0A5E}].
\end{proof}

% object: remark-properness-needed-c7 (blueprint/chap-Proper-base-change/sec-The-projective-line/subsec-Henselian-pairs-and-curves/remark-properness-needed-c7/_section.md)
\phantomsection\label{bp62}%
\begin{bpstatement}{Remark~4.22}{ (Why properness is needed in $(C_7)$)}
Properness in \((C_7)\) is needed when the torsion order is divisible by the characteristic, and this case arises in Proposition \hyperref[bp76]{4.26}, where the field is extended for arbitrary torsion sheaves. For example, in characteristic \(p\),

\[
H^1(\mathbf A^1_k,\mathbf F_p)
=k[t]/\{h^p-h:h\in k[t]\},
\]

and extending an algebraically closed ground field adds new classes. Indeed, we have \(c^pt^{pm}\equiv ct^m\), and every element of \(k\) is of the form \(h^p-h\). It follows that every class has a unique representative \(\sum_{p\nmid m}a_mt^m\), and the same holds over an algebraically closed \(k'\supset k\). Hence for \(a\in k'\smallsetminus k\) the class of \(at\) does not come from \(k[t]\).
\end{bpstatement}


% object: subsec-Injectives-on-the-closed-fibre (blueprint/chap-Proper-base-change/sec-The-projective-line/subsec-Injectives-on-the-closed-fibre/_section.md), content hidden
\subsubsection*{\texorpdfstring{4.4.2\quad Injectives on the closed fibre}{4.4.2 Injectives on the closed fibre}}
\phantomsection\label{bp63}%
\addcontentsline{toc}{subsubsection}{\texorpdfstring{4.4.2\quad Injectives on the closed fibre}{4.4.2 Injectives on the closed fibre}}
% text: prose-04-24
\phantomsection\label{bp64}%
Let \(R\) be a henselian local ring with maximal ideal \(\mathfrak m\) and residue field \(\kappa\), and let \(i:\mathbf P^1_\kappa\hookrightarrow\mathbf P^1_R\) be the inclusion of the closed fibre. For a sheaf \(G\) on \(\mathbf P^1_{R,\acute{e}t}\) write \(G_0=i^*G\). To compare the cohomology of \(\mathbf P^1_R\) with that of its closed fibre, we apply the comparison in degree zero of Theorem \hyperref[bp24]{4.8} to an injective resolution. This works once we know that the restriction \(I_0\) of an injective sheaf \(I\) is acyclic on \(\mathbf P^1_\kappa\) (Lemma \hyperref[bp68]{4.24}). However, an injective sheaf is too large to work with directly, so we approximate it by constructible sheaves and embed these into direct images along finite morphisms, using three standard finiteness facts about small étale topoi of quasi-compact quasi-separated schemes.

% object: lemma-finiteness-small-etale-topoi (blueprint/chap-Proper-base-change/sec-The-projective-line/subsec-Injectives-on-the-closed-fibre/lemma-finiteness-small-etale-topoi/_section.md)
\phantomsection\label{bp65}%
\begin{bpstatementit}{Lemma~4.23}{ (Finiteness facts on $\mathbf P^1_R$)}
\begin{enumerate}
\item A sheaf of \(\mathbf F_\ell\)-modules on \(\mathbf P^1_{R,\acute{e}t}\) is a filtered colimit of constructible ones, and cohomology on \(\mathbf P^1_R\) and on \(\mathbf P^1_\kappa\) commutes with filtered colimits.
\item A constructible sheaf \(G\) of \(\mathbf F_\ell\)-modules on \(\mathbf P^1_{R,\acute{e}t}\) admits a monomorphism \[
   G\hookrightarrow
   \bigoplus_{a=1}^N(h_a)_*\underline{M_a},
   \qquad h_a:Y_a\to \mathbf P^1_R\ \text{finite and finitely presented},
   \tag{4.9}
   \] with \(M_a\) finite \(\mathbf F_\ell\)-modules.
\item For a finite morphism \(h:Y\to \mathbf P^1_R\) and an abelian sheaf \(G\) on \(Y_{\acute{e}t}\), we have \[
   R^qh_*G=0\quad(q>0),\qquad
   i^*h_*G\xrightarrow{\sim}(h_0)_*i_Y^*G.
   \tag{4.10}
   \] Here \(h_0:Y_0=Y\times_{\mathbf P^1_R}\mathbf P^1_\kappa\to \mathbf P^1_\kappa\) and \(i_Y:Y_0\to Y\) are the base changes of \(h\) and \(i\), and the second map is the underived base change map (4.2).
\end{enumerate}
\end{bpstatementit}

% proof: proof/lemma-finiteness-small-etale-topoi
\begin{proof}[{Proof}]\phantomsection\label{bp66}%
Part (1) is [Stacks, \href{https://stacks.math.columbia.edu/tag/03SA}{03SA}] and [Stacks, \href{https://stacks.math.columbia.edu/tag/03Q5}{03Q5}], and part (2) is [Stacks, \href{https://stacks.math.columbia.edu/tag/09Z7}{09Z7}]. Part (3) is the case \(u=h\), \(W=\mathbf P^1_R\) of Lemma \hyperref[bp80]{4.27} below, which is proved directly from the stalk description of finite direct images [Stacks, \href{https://stacks.math.columbia.edu/tag/03QP}{03QP}].
\end{proof}

% text: prose-04-25
\phantomsection\label{bp67}%
The sheaves on the right of (4.9) are direct images along finite morphisms, and by Lemma lemma-finiteness-small-etale-topoi such direct images are exact functors and commute with restriction to the closed fibre.

% object: lemma-acyclicity-closed-fibre (blueprint/chap-Proper-base-change/sec-The-projective-line/subsec-Injectives-on-the-closed-fibre/lemma-acyclicity-closed-fibre/_section.md)
\phantomsection\label{bp68}%
\begin{bpstatementit}{Lemma~4.24}{ (Acyclicity on the closed fibre)}
If \(I\) is an injective sheaf of \(\mathbf F_\ell\)-modules on \(\mathbf P^1_{R,\acute{e}t}\), then \(H^q(\mathbf P^1_\kappa,I_0)=0\) for every \(q>0\).
\end{bpstatementit}

% proof: proof/lemma-acyclicity-closed-fibre
\begin{proof}[{Proof}]\phantomsection\label{bp69}%
We prove vanishing on affine charts by Gabber's theorem, deduce the case \(q\geq2\) by Mayer–Vietoris, and treat \(q=1\) by lifting classes to finite torsors over \(\mathbf P^1_R\), where injectivity of \(I\) kills them.

\emph{Step 1: affine opens.} Let \(W=\operatorname{Spec}B\) be an affine open of \(\mathbf P^1_R\) and set \(J=\mathfrak mB\). We show that \(I_0|_{W_0}\) has no cohomology in positive degrees. Since the pair \((B,J)\) need not be henselian, we pass to its henselization \((B^h,JB^h)\), for which \(B^h/JB^h=B/J\). Moreover, \(B^h\) is the filtered colimit of the étale \(B\)-algebras \(B'\) equipped with \(B'/JB'\simeq B/J\). Gabber's theorem (Theorem \hyperref[bp51]{4.18}) and continuity of étale cohomology give

\[
H^q(W_0,I_0|_{W_0})
\simeq H^q(\operatorname{Spec}B^h,I|)
\simeq\varinjlim_{B'}H^q(\operatorname{Spec}B',I|).
\tag{4.11}
\]

Every group on the right vanishes for \(q>0\), because the restriction of an injective sheaf to an étale object is injective: the restriction functor has an exact left adjoint, namely extension by zero in the localized site.

\emph{Step 2: degrees \(q\geq2\).} The two standard affine charts \(U,V\subset \mathbf P^1_R\) have affine intersection, so (4.11) shows that the restriction of \(I\) to each of \(U_0\), \(V_0\) and \((U\cap V)_0\) has no cohomology in positive degrees. The Mayer–Vietoris sequence on \(\mathbf P^1_\kappa\) therefore gives \(H^q(\mathbf P^1_\kappa,I_0)=0\) for \(q\geq2\) (this is the case of two charts in [Stacks, \href{https://stacks.math.columbia.edu/tag/0A51}{0A51}]). In degree one the same sequence only identifies \(H^1(\mathbf P^1_\kappa,I_0)\) with the cokernel of \(\Gamma(U_0,I_0)\oplus\Gamma(V_0,I_0)\to\Gamma((U\cap V)_0,I_0)\), and nothing said so far shows that this cokernel vanishes.

\emph{Step 3: degree one.} Let \(\xi\in H^1(\mathbf P^1_\kappa,I_0)\). Following [Stacks, \href{https://stacks.math.columbia.edu/tag/0A5H}{0A5H}], we show that \(\xi=0\). First write \(I\) as a filtered colimit of constructible \(\mathbf F_\ell\)-sheaves. Since pullback and \(H^1\) on \(\mathbf P^1_\kappa\) commute with this colimit, \(\xi\) is the image of a class \(\zeta\in H^1(\mathbf P^1_\kappa,G_0)\) for some constructible \(G\) with a map \(G\to I\). Next choose (4.9) and put \(P=\bigoplus_a(h_a)_*\underline{M_a}\). Since \(I\) is injective, the map \(G\to I\) extends to \(P\to I\). Let \(\zeta'\) be the image of \(\zeta\) in \(H^1(\mathbf P^1_\kappa,P_0)\).

The class \(\zeta'\) lifts to \(H^1(\mathbf P^1_R,P)\). Indeed, (4.10) and the Leray spectral sequence give

\[
H^1(\mathbf P^1_\kappa,P_0)
\simeq\bigoplus_a H^1((Y_a)_0,\underline{M_a}),\qquad
H^1(\mathbf P^1_R,P)\simeq\bigoplus_a H^1(Y_a,\underline{M_a}).
\]

Each \(Y_a\) is finite over \(\mathbf P^1_R\), hence proper over \(R\), so the comparison (4.8) for finite torsors (Theorem \hyperref[bp53]{4.19}) lifts \(\zeta'\). This lift maps to zero in \(H^1(\mathbf P^1_R,I)\) because \(I\) is injective. Now the commutative square of restriction maps

\[
\begin{CD}
H^1(\mathbf P^1_R,P) @>>> H^1(\mathbf P^1_R,I)\\
@VVV @VVV\\
H^1(\mathbf P^1_\kappa,P_0) @>>> H^1(\mathbf P^1_\kappa,I_0)
\end{CD}
\]

shows that the image of \(\zeta'\) in \(H^1(\mathbf P^1_\kappa,I_0)\) is zero, and this image is \(\xi\), because \(G\to P\to I\) is the given map \(G\to I\). Hence \(\xi=0\).
\end{proof}

% text: prose-04-26
\phantomsection\label{bp70}%
With this acyclicity, an injective resolution on \(\mathbf P^1_R\) restricts to an acyclic resolution on the closed fibre. Applying the degree-zero comparison of Theorem \hyperref[bp24]{4.8} term by term, we obtain the comparison in all degrees. Since the lemma only concerns sheaves killed by a prime \(\ell\), we then extend the result to all torsion sheaves by dévissage.

% object: proposition-cohomology-p-1-over-henselian (blueprint/chap-Proper-base-change/sec-The-projective-line/subsec-Injectives-on-the-closed-fibre/proposition-cohomology-p-1-over-henselian/_section.md)
\phantomsection\label{bp71}%
\begin{bpstatementit}{Proposition~4.25}{ (Cohomology of $\mathbf P^1$ over henselian rings)}
Let \(R\) be a henselian local ring with residue field \(\kappa\). For every torsion abelian sheaf \(F\) on \(\mathbf P^1_{R,\acute{e}t}\), restriction induces isomorphisms, natural in \(F\), with \(F_0=i^*F\),

\[
H^q(\mathbf P^1_R,F)\xrightarrow{\sim}
H^q(\mathbf P^1_\kappa,F_0)\qquad(q\geq0).
\]
\end{bpstatementit}

% proof: proof/proposition-cohomology-p-1-over-henselian
\begin{proof}[{Proof}]\phantomsection\label{bp72}%
We treat sheaves killed by a prime first, where Lemma \hyperref[bp68]{4.24} applies, and then reach all torsion sheaves by extensions, direct sums and filtered colimits.

Suppose first that \(\ell F=0\) for a prime \(\ell\).  Take an injective resolution \(F\to I^\bullet\) in sheaves of \(\mathbf F_\ell\)-modules on \(\mathbf P^1_{R,\acute{e}t}\). Since \(i^*\) is exact, \(F_0\to I_0^\bullet\) is again a resolution, and by Lemma \hyperref[bp68]{4.24} its terms are \(\Gamma(\mathbf P^1_\kappa,-)\)-acyclic. Moreover, Theorem \hyperref[bp24]{4.8} shows that the two complexes of sections agree term by term:

\[
\Gamma(\mathbf P^1_R,I^\bullet)\xrightarrow{\sim}
\Gamma(\mathbf P^1_\kappa,I_0^\bullet).
\]

The left complex computes \(R\Gamma(\mathbf P^1_R,F)\), and the right computes \(R\Gamma(\mathbf P^1_\kappa,F_0)\). Hence the natural restriction map between them is a quasi-isomorphism.

By the long exact cohomology sequences and the five lemma, the sheaves for which the statement holds form a class closed under extensions. Since a sheaf killed by \(\ell^a\) has a finite filtration with quotients killed by \(\ell\), and a sheaf killed by an integer \(n\) is the finite direct sum of its primary parts, the class contains all sheaves of bounded torsion. An arbitrary torsion \(F\) is the colimit \(\varinjlim_n F[n]\), with positive integers ordered by divisibility, and the comparison passes to this colimit, because the cohomology of the quasi-compact quasi-separated schemes \(\mathbf P^1_R\) and \(\mathbf P^1_\kappa\) commutes with filtered colimits [Stacks, \href{https://stacks.math.columbia.edu/tag/03Q5}{03Q5}]. Finally, all maps used are restriction maps, which gives naturality in \(F\).
\end{proof}

% text: prose-04-27
\phantomsection\label{bp73}%
For example, let \(R\) be strictly henselian and \(n\) invertible in \(R\). Then the residue field \(\kappa\) is separably closed, and passing to its algebraic closure does not change the étale topos [Stacks, \href{https://stacks.math.columbia.edu/tag/03SI}{03SI}]. Hence \(H^2(\mathbf P^1_\kappa,\mu_n)\simeq\mathbf Z/n\), generated by the Kummer class of \(\mathcal O(1)\) (Proposition \bpout{D.2}, in Appendix \bpout{D}), and the proposition says that \(H^2(\mathbf P^1_R,\mu_n)\) is also \(\mathbf Z/n\), generated by the class of \(\mathcal O(1)\).


% object: subsec-Arbitrary-base-change (blueprint/chap-Proper-base-change/sec-The-projective-line/subsec-Arbitrary-base-change/_section.md), content hidden
\subsubsection*{\texorpdfstring{4.4.3\quad Arbitrary base change}{4.4.3 Arbitrary base change}}
\phantomsection\label{bp74}%
\addcontentsline{toc}{subsubsection}{\texorpdfstring{4.4.3\quad Arbitrary base change}{4.4.3 Arbitrary base change}}
% text: prose-04-28
\phantomsection\label{bp75}%
Let \(\pi:\mathbf P^1_S\to S\) be the projection and \(\bar s\) a geometric point of \(S\). The stalk of \(R^q\pi_*F\) at \(\bar s\) is the cohomology of \(\mathbf P^1\) over the strict localization at \(\bar s\), by (4.12) below. By Proposition \hyperref[bp71]{4.25} it is therefore the cohomology of \(\mathbf P^1_{k_s}\), where \(k_s\) is the separably closed residue field. Along a base change \(g:T\to S\) the residue field grows from \(k_s\) to \(k_t\), so on stalks the base change map becomes the pullback from \(\mathbf P^1_{k_s}\) to \(\mathbf P^1_{k_t}\), which is an isomorphism by Lemma \hyperref[bp60]{4.21}. This is the second place where properness enters: for \(\mathbf A^1\) and \(p\)-torsion coefficients the corresponding map is not surjective.

The stalk formula holds for every quasi-compact quasi-separated morphism \(f:Z\to S\). For a geometric point \(\bar s\) of \(S\) and an abelian sheaf \(G\) on \(Z_{\acute{e}t}\) we have

\[
(R^qf_*G)_{\bar s}
\simeq H^q(Z\times_SS_{(\bar s)},G|),\qquad
S_{(\bar s)}=\operatorname{Spec}\mathcal O^{\mathrm{sh}}_{S,\bar s}.
\tag{4.12}
\]

Indeed, the stalk is the filtered colimit of \(H^q(Z\times_SV,G|)\) over pointed étale neighbourhoods \((V,\bar v)\) of \(\bar s\), and by continuity of étale cohomology this colimit is the cohomology of the indicated limit scheme [Stacks, \href{https://stacks.math.columbia.edu/tag/03Q9}{03Q9}].

% object: proposition-proper-base-change-projective-line (blueprint/chap-Proper-base-change/sec-The-projective-line/subsec-Arbitrary-base-change/proposition-proper-base-change-projective-line/_section.md)
\phantomsection\label{bp76}%
\begin{bpstatementit}{Proposition~4.26}{ (Proper base change for $\mathbf P^1$)}
Let \(S\) be a scheme and \(\Lambda\) a commutative ring.  In the cartesian square

\[
\begin{CD}
\mathbf P^1_T @>{g'}>> \mathbf P^1_S\\
@V{\pi'}VV @VV{\pi}V\\
T @>{g}>> S,
\end{CD}
\]

the base change map

\[
g^*R\pi_*K\longrightarrow R\pi'_*g'^*K
\]

is an isomorphism for every \(K\in D^+(\mathbf P^1_{S,\acute{e}t},\Lambda)\) with torsion cohomology sheaves. In particular, \(\mathrm{IA}_\Lambda(\pi)\) holds.
\end{bpstatementit}

% proof: proof/proposition-proper-base-change-projective-line
\begin{proof}[{Proof}]\phantomsection\label{bp77}%
For a torsion sheaf we compare the two sides on geometric stalks. Both stalks are the cohomology of a projective line over a separably closed field, and the base change map becomes the pullback along a field extension. Bounded below complexes then follow by truncation.

Let \(F\) be a torsion sheaf. Choose a geometric point \(\bar t:\operatorname{Spec}\Omega\to T\) and put \(\bar s=g\bar t\). Then the strict localizations fit into a natural map \(T_{(\bar t)}\to S_{(\bar s)}\) of strictly henselian local schemes. Write \(k_s,k_t\) for their residue fields; these are separably closed, and the map induces an inclusion \(k_s\hookrightarrow k_t\). For each \(q\geq0\), using (4.12) and Proposition \hyperref[bp71]{4.25}, we obtain the commutative diagram

\[
\begin{CD}
(g^*R^q\pi_*F)_{\bar t}
@>>>
(R^q\pi'_*g'^*F)_{\bar t}\\
@V{\sim}VV @VV{\sim}V\\
H^q(\mathbf P^1_{S_{(\bar s)}},F|)
@>>> H^q(\mathbf P^1_{T_{(\bar t)}},F|)\\
@V{\sim}VV @VV{\sim}V\\
H^q(\mathbf P^1_{k_s},F|)
@>>> H^q(\mathbf P^1_{k_t},F|).
\end{CD}
\]

Here the top arrow is the stalk at \(\bar t\) of the map induced by (4.3) on \(q\)th cohomology sheaves, while the middle and bottom arrows are the natural pullback maps. The top square commutes by the definition of the base change map: a section over \(\mathbf P^1_V\), for a pointed étale neighbourhood \(V\to S\), is pulled back to \(\mathbf P^1_{V\times_ST}\), and passing to the limits over the two strict localizations gives exactly the middle arrow. The bottom square commutes by naturality of restriction to the closed fibres.

The vertical arrows are isomorphisms: the upper ones by (4.12), applied to \(\pi\) and to \(\pi'\); here \(\mathbf P^1_{S_{(\bar s)}}\) and \(\mathbf P^1_{T_{(\bar t)}}\) are the base changes of \(\mathbf P^1_S\) and \(\mathbf P^1_T\) to the strict localizations. The lower ones are isomorphisms by Proposition \hyperref[bp71]{4.25}, applied to the strictly henselian local rings of \(S_{(\bar s)}\) and \(T_{(\bar t)}\). Hence the top arrow is an isomorphism as soon as the bottom arrow is, and the bottom arrow is an isomorphism by Lemma \hyperref[bp60]{4.21}, applied to \(C=\mathbf P^1_{k_s}\) and the extension \(k_s\subseteq k_t\). Thus every geometric stalk of \(g^*R^q\pi_*F\to R^q\pi'_*g'^*F\) is an isomorphism. As geometric stalks are conservative, the base change map for \(F\) is an isomorphism in \(D^+(T_{\acute{e}t},\mathbf Z)\). Moreover, cohomology and higher direct images of sheaves of \(\Lambda\)-modules agree with those of their underlying abelian sheaves (Lemma \hyperref[bp8]{4.2}). Hence the base change map for \(F\) is also an isomorphism in \(D^+(T_{\acute{e}t},\Lambda)\).

Now let \(K\in D^+(\mathbf P^1_{S,\acute{e}t},\Lambda)\) have torsion cohomology and fix \(m\in\mathbf Z\). The two functors \(g^*R\pi_*,\ R\pi'_*g'^*:D^+(\mathbf P^1_{S,\acute{e}t},\Lambda)\to
D^+(T_{\acute{e}t},\Lambda)\) send \(D^{\geq m+1}\) into \(D^{\geq m+1}\), since inverse image is exact and derived direct image is left \(t\)-exact. Using the truncation triangle for \(\tau_{\leq m}K\), we can therefore reduce the assertion in degree \(m\) to the bounded complex \(\tau_{\leq m}K\). By induction on the length, using the truncation triangles and naturality of the base change map, we reduce further to the shifted torsion sheaves \(H^a(K)[-a]\). For these the assertion has already been proved.

For \(\mathrm{IA}_\Lambda(\pi)\), let \(I\) be a torsion injective sheaf of \(\Lambda\)-modules. Then \(R^q\pi_*I=0\) for \(q>0\), and the case of sheaves gives \(R^q\pi'_*g'^*I\simeq g^*R^q\pi_*I=0\).
\end{proof}



% object: sec-Proper-morphisms (blueprint/chap-Proper-base-change/sec-Proper-morphisms/_section.md), content hidden
\subsection*{\texorpdfstring{4.5\quad Proper morphisms}{4.5 Proper morphisms}}
\phantomsection\label{bp78}%
\addcontentsline{toc}{subsection}{\texorpdfstring{4.5\quad Proper morphisms}{4.5 Proper morphisms}}
% text: prose-04-29
\phantomsection\label{bp79}%
Proper base change for an arbitrary proper morphism follows from the case of the projective line (Proposition \hyperref[bp76]{4.26}) by formal steps. Finite morphisms satisfy proper base change directly, and Lemma \hyperref[bp85]{4.29} shows that the property is stable under composition and descends along proper surjections. Descent along a proper surjection \(u:X\to Y\) works because every injective sheaf on \(Y\) is a direct summand of \(u_*J\) for an injective sheaf \(J\) on \(X\), and the hypotheses on \(u\) and on the composite keep \(u_*J\) acyclic after base change.

The steps are then as follows. First, composition gives the projection of \((\mathbf P^1)^n\). Then descent along the finite surjection \((\mathbf P^1)^n\to\mathbf P^n\) gives the projection of \(\mathbf P^n\), and with it every projective morphism. Finally, by Chow's lemma every proper morphism is dominated by a projective one through a proper surjection, so a last descent covers all proper morphisms.

% object: lemma-finite-morphisms (blueprint/chap-Proper-base-change/sec-Proper-morphisms/lemma-finite-morphisms/_section.md)
\phantomsection\label{bp80}%
\begin{bpstatementit}{Lemma~4.27}{ (Finite morphisms)}
Let \(u:V\to W\) be finite. For an abelian sheaf \(F\) and a geometric point \(\bar w:\operatorname{Spec}(\Omega)\to W\), with \(\Omega\) algebraically closed, one has

\[
(u_*F)_{\bar w}
\simeq \prod_{\substack{\bar v:\operatorname{Spec}(\Omega)\to V\\
u\bar v=\bar w}}F_{\bar v}.
\]

The product is finite. The functor \(u_*\) from abelian sheaves on \(V_{\acute{e}t}\) to abelian sheaves on \(W_{\acute{e}t}\) is exact, and its formation commutes with arbitrary base change. Consequently, for every base change \(W'\to W\) and every constant coefficient ring \(\Lambda\), the derived base change map (4.3) is an isomorphism for every \(K\in D^+(V_{\acute{e}t},\Lambda)\).
\end{bpstatementit}

% proof: proof/lemma-finite-morphisms
\begin{proof}[{Proof}]\phantomsection\label{bp81}%
Everything follows from the stalk formula. To prove it, note that over a strict localization the source splits into strictly henselian local pieces, and on such a piece taking sections is the same as taking a stalk.

After base change to the strict henselization of \(W\) at \(\bar w\), the finite scheme \(V\) becomes a finite disjoint union of strictly henselian local schemes. On each such local scheme, the sections of an étale sheaf are its stalk at the closed geometric point. Indeed, every étale neighbourhood of that point admits a section, so a germ extends to a global section, and two sections with the same germ agree after pulling back along such a section. In particular, taking sections is exact on abelian sheaves. Combining this with the description of stalks of direct images via strict localizations [Stacks, \href{https://stacks.math.columbia.edu/tag/03QP}{03QP}], we obtain the displayed formula. Since geometric stalks are conservative and finite products of abelian groups are exact, \(u_*\) is exact.

In a cartesian square over \(W'\to W\), a geometric point of \(W'\) and its image in \(W\) have the same geometric fibre. Therefore the underived base change map is the identity on the corresponding factors of the displayed formula, hence an isomorphism. Since both direct images in the square are exact, their derived functors and the base change map are computed term by term. Finally, the same argument applies to sheaves of modules.
\end{proof}

% text: prose-04-30
\phantomsection\label{bp82}%
For instance, if \(u:\operatorname{Spec}k'\to\operatorname{Spec}k\) is a separable extension of degree \(d\) and \(\bar w\) is a geometric point, the formula says that \((u_*F)_{\bar w}\) is the product of the \(d\) stalks of \(F\) at the \(d\) embeddings of \(k'\) into \(\Omega\) over \(k\).

% object: example-finite-direct-image-double-cover (blueprint/chap-Proper-base-change/sec-Proper-morphisms/example-finite-direct-image-double-cover/_section.md)
\phantomsection\label{bp83}%
\begin{bpstatement}{Example~4.28}{ (A ramified double cover)}
Let \(k\) be algebraically closed of characteristic \(\neq2\), let \(u:\mathbf A^1_k\to\mathbf A^1_k\) be \(x\mapsto x^2\), and let \(A\) be a finite ring. Over a point \(t\neq0\) the fibre consists of the two points \(\pm\sqrt t\), whereas over \(t=0\) it is the single point \(x=0\), with a nonreduced structure that the étale topology does not see. By the lemma, therefore,

\[
(u_*A)_{\bar t}=A\oplus A\quad(t\neq0),\qquad (u_*A)_{\bar0}=A,
\]

and the rank of the stalks jumps at the origin, so \(u_*A\) is not lisse. It is, however, constructible for the stratification \(\mathbf G_m\sqcup\{0\}\). Base change to the origin gives the stalk \(A\), the cohomology of the fibre, as the lemma predicts. For the open immersion \(j:\mathbf G_m\to\mathbf A^1\), by contrast, the stalk of \(j_*A\) at \(0\) is \(A\) while the fibre is empty.
\end{bpstatement}

% text: prose-04-31
\phantomsection\label{bp84}%
Stability under composition and descent along proper surjections are best proved on the level of injective acyclicity, since by Theorem \hyperref[bp45]{4.17} it is equivalent to proper base change and it lends itself to a spectral sequence argument.

% object: lemma-composition-cancellation (blueprint/chap-Proper-base-change/sec-Proper-morphisms/lemma-composition-cancellation/_section.md)
\phantomsection\label{bp85}%
\begin{bpstatementit}{Lemma~4.29}{ (Composition and cancellation)}
Let \(X\xrightarrow{u}Y\xrightarrow{v}Z\) be proper morphisms. Then proper base change is stable under composition, and it descends along a proper surjection:

\begin{itemize}
\item (i) If \(\mathrm{PBC}_\Lambda(u)\) and \(\mathrm{PBC}_\Lambda(v)\) hold, then \(\mathrm{PBC}_\Lambda(vu)\) holds.
\item (ii) If \(u\) is surjective and \(\mathrm{PBC}_\Lambda(u)\) and \(\mathrm{PBC}_\Lambda(vu)\) hold, then \(\mathrm{PBC}_\Lambda(v)\) holds.
\end{itemize}
\end{bpstatementit}

% proof: proof/lemma-composition-cancellation
\begin{proof}[{Proof}]\phantomsection\label{bp86}%
Theorem \hyperref[bp45]{4.17} reduces both parts to the condition \(\mathrm{IA}_\Lambda\) of injective acyclicity, and we read off this vanishing from the Leray spectral sequence of the base-changed composite. For (ii) the hypotheses do not apply directly to an injective sheaf on \(Y\), so we first realize it as a direct summand of \(u_*J\) with \(J\) injective on \(X\).

To set up the spectral sequence, fix an arbitrary base change \(Z'\to Z\). Write \(g_X:X'\to X\) and \(g_Y:Y'\to Y\) for the projections, and \(u':X'\to Y'\), \(v':Y'\to Z'\) for the induced maps. The direct image \(u'_*\) preserves injectives, because its left adjoint \(u'^*\) on sheaves of \(\Lambda\)-modules is exact. Hence for every sheaf \(J\) of \(\Lambda\)-modules on \(X\) we have the Grothendieck spectral sequence for the composite direct image [Stacks, \href{https://stacks.math.columbia.edu/tag/0734}{0734}]:

\[
R^pv'_*R^qu'_*g_X^*J
\ \Longrightarrow\ R^{p+q}(v'u')_*g_X^*J.
\tag{4.13}
\]

\emph{Part (i).} Let \(J\) be an injective sheaf of \(\Lambda\)-modules on \(X\). We show \(R^m(v'u')_*g_X^*J=0\) for \(m>0\), which is \(\mathrm{IA}_\Lambda(vu)\). For this it suffices that the \(E_2\)-terms of (4.13) vanish except for \(p=q=0\), and they do. Indeed, \(R^qu'_*g_X^*J=0\) for \(q>0\) by \(\mathrm{IA}_\Lambda(u)\). For \(q=0\), proper base change in degree zero (Proposition \hyperref[bp30]{4.11}) gives

\[
u'_*g_X^*J\simeq g_Y^*u_*J,
\]

and the sheaf \(u_*J\) is injective, so \(R^pv'_*g_Y^*u_*J=0\) for \(p>0\) by \(\mathrm{IA}_\Lambda(v)\). The abutment therefore vanishes in positive degrees.

\emph{Part (ii), the splitting.} Let \(I\) be an injective sheaf of \(\Lambda\)-modules on \(Y\). Choose a monomorphism of sheaves

\[
u^*I\hookrightarrow J
\]

with \(J\) injective on \(X\), and let \(\iota:I\to u_*J\) be its adjoint. The map \(\iota\) is a monomorphism. To see this, given a geometric point \(\bar y\), choose a geometric point \(\bar x\) above it, which exists because \(u\) is surjective. Then the composite

\[
I_{\bar y}\xrightarrow{\iota_{\bar y}}(u_*J)_{\bar y}
\longrightarrow J_{\bar x}
\]

is the stalk of the chosen monomorphism, under \((u^*I)_{\bar x}=I_{\bar y}\), so it is injective. As \(I\) is an injective sheaf, there is then a map of sheaves \(s:u_*J\to I\) with \(s\iota=\mathrm{id}_I\). In other words, \(I\) is a direct summand of \(u_*J\).

\emph{Part (ii), the vanishing.} We first show \(R^pv'_*g_Y^*u_*J=0\) for \(p>0\). The hypotheses \(\mathrm{IA}_\Lambda(u)\) and \(\mathrm{IA}_\Lambda(vu)\) give

\[
R^qu'_*g_X^*J=0\quad(q>0),\qquad
R^m(v'u')_*g_X^*J=0\quad(m>0),
\]

so the spectral sequence (4.13) collapses to its row \(q=0\). This means that \(R^pv'_*u'_*g_X^*J\simeq R^p(v'u')_*g_X^*J=0\) for \(p>0\). Since \(u'_*g_X^*J\simeq g_Y^*u_*J\) by proper base change in degree zero, we obtain

\[
R^pv'_*g_Y^*u_*J
\simeq R^pv'_*u'_*g_X^*J=0\qquad(p>0).
\]

Now \(g_Y^*I\) is a direct summand of \(g_Y^*u_*J\), by applying \(g_Y^*\) to the splitting \(s\iota=\mathrm{id}_I\), and as \(R^pv'_*\) is additive, \(R^pv'_*g_Y^*I\) is a direct summand of \(R^pv'_*g_Y^*u_*J\). Hence \(R^pv'_*g_Y^*I=0\) for \(p>0\), which is \(\mathrm{IA}_\Lambda(v)\).
\end{proof}

% text: prose-04-32
\phantomsection\label{bp87}%
Part (i) already gives proper base change for \((\mathbf P^1_S)^n\to S\), a composite of projections of projective lines. To reach \(\mathbf P^n_S\) we use part (ii), and for this we need a finite surjection from \((\mathbf P^1_S)^n\) onto \(\mathbf P^n_S\) over \(S\).

% object: lemma-finite-cover-projective-space (blueprint/chap-Proper-base-change/sec-Proper-morphisms/lemma-finite-cover-projective-space/_section.md)
\phantomsection\label{bp88}%
\begin{bpstatementit}{Lemma~4.30}{ (A finite cover of projective space)}
For every scheme \(S\) and every \(n\geq1\) there is a finite surjective \(S\)-morphism

\[
\sigma_n:(\mathbf P^1_S)^n\longrightarrow\mathbf P^n_S.
\]
\end{bpstatementit}

% proof: proof/lemma-finite-cover-projective-space
\begin{proof}[{Proof}]\phantomsection\label{bp89}%
Recall that \(\mathbf P^n\) parametrizes binary forms of degree \(n\) up to scalars, and the map \(\sigma_n\) multiplies \(n\) linear forms; over a field it is the quotient of \((\mathbf P^1)^n\) by the symmetric group. Explicitly, we use the coefficients of a product of binary linear forms:

\[
\prod_{i=1}^n(x_iT+y_iU)
=\sum_{j=0}^nF_jT^{n-j}U^j,\qquad
\sigma_n\bigl(([x_i:y_i])_i\bigr)=[F_0:\cdots:F_n].
\]

Every \(F_j\) is homogeneous of degree one in each pair \((x_i,y_i)\), and the corresponding sections of \(\mathcal O(1,\ldots,1)\) have no common zero. Indeed, at a geometric point every linear factor is nonzero, and hence so is their product. Therefore the sections define the stated morphism over any \(S\), and the construction commutes with base change.

For example, on the chart \(F_0\ne0\) of the target, whose preimage is the chart where every \(x_i\) is invertible, the coordinates \(c_j=F_j/F_0\) pull back to the elementary symmetric polynomials in the \(w_i=y_i/x_i\), as in the classical description of \(\mathbf A^n\) as the quotient of \(\mathbf A^n\) by the symmetric group.

Finiteness of \(\sigma_n\) follows from properness and finiteness of the fibres. Since the source is proper over \(S\) and the target is separated over \(S\), the morphism \(\sigma_n\) is proper. For the fibres, recall that over an algebraically closed field, a nonzero binary form of degree \(n\) is a product of \(n\) linear forms \(x_iT+y_iU\), and these are unique up to order and up to multiplying each by a nonzero constant. Here a factor with \(x_i=0\) is a multiple of \(U\) and corresponds to the point \([0:1]\). In other words, the form determines the multiset of points \([x_i:y_i]\in\mathbf P^1\). Hence every geometric fibre of \(\sigma_n\) is nonempty, and \(\sigma_n\) is surjective. Moreover, each geometric fibre has at most \(n!\) geometric points, namely the orderings of this multiset. Since a scheme of finite type over a field with finitely many geometric points is zero dimensional, \(\sigma_n\) is quasi-finite, and a proper quasi-finite morphism is finite [Stacks, \href{https://stacks.math.columbia.edu/tag/02LS}{02LS}].
\end{proof}

% text: prose-04-33
\phantomsection\label{bp90}%
For \(n=1\) the map \(\sigma_1\) is the identity of \(\mathbf P^1_S\). For \(n=2\) it sends \(([x_1:y_1],[x_2:y_2])\) to \([x_1x_2:x_1y_2+y_1x_2:y_1y_2]\), and over a field it is the quotient by swapping the two factors, of degree \(2\).

% object: lemma-pbc-projective-space (blueprint/chap-Proper-base-change/sec-Proper-morphisms/lemma-pbc-projective-space/_section.md)
\phantomsection\label{bp91}%
\begin{bpstatementit}{Lemma~4.31}{ (Projective space)}
If \(\mathrm{PBC}_\Lambda(\mathbf P^1_S\to S)\) holds for every scheme \(S\), then \(\mathrm{PBC}_\Lambda(\mathbf P^n_S\to S)\) holds for every scheme \(S\) and every \(n\geq0\).
\end{bpstatementit}

% proof: proof/lemma-pbc-projective-space
\begin{proof}[{Proof}]\phantomsection\label{bp92}%
The case \(n=0\) is the identity. For \(n\geq1\), the projection \((\mathbf P^1_S)^n\to S\) is a composite of \(n\) projections of projective lines over successive bases. Hence it satisfies \(\mathrm{PBC}_\Lambda\) by the hypothesis and Lemma lemma-composition-cancellation. The finite surjection \(\sigma_n\) of Lemma \hyperref[bp88]{4.30} is a morphism over \(S\), and it satisfies \(\mathrm{PBC}_\Lambda\) by Lemma \hyperref[bp80]{4.27}. Cancelling \(\sigma_n\) by Lemma lemma-composition-cancellation, we obtain \(\mathrm{PBC}_\Lambda\) for \(\mathbf P^n_S\to S\).
\end{proof}

% text: prose-04-34
\phantomsection\label{bp93}%
We can now treat projective morphisms, since they factor as a closed immersion followed by the projection of a projective space. By Chow's lemma an arbitrary proper morphism is dominated by a projective one through a proper surjection, and the cancellation of Lemma lemma-composition-cancellation then passes proper base change to the proper morphism.

% object: proposition-proper-base-change-etale-site (blueprint/chap-Proper-base-change/sec-Proper-morphisms/proposition-proper-base-change-etale-site/_section.md)
\phantomsection\label{bp94}%
\begin{bpstatementit}{Proposition~4.32}{ (Étale proper base change)}
If \(\mathrm{PBC}_\Lambda(\mathbf P^1_S\to S)\) holds for every scheme \(S\), then \(\mathrm{PBC}_\Lambda(f)\) holds for every proper morphism of schemes \(f\).
\end{bpstatementit}

% proof: proof/proposition-proper-base-change-etale-site
\begin{proof}[{Proof}]\phantomsection\label{bp95}%
Let \(f:X\to Y\) be proper. We first reduce to the case that \(Y\) is affine, as Chow's lemma requires. Let \(Y=\bigcup U\) be an affine open cover, and assume \(\mathrm{PBC}_\Lambda\) for each restriction \(f^{-1}(U)\to U\). Given a base change \(Y'\to Y\), the inverse images \(U'\) of the sets \(U\) then cover \(Y'\). Since higher direct images commute with restriction to open subschemes, the restriction of the base change map to \(U'\) is the base change map of \(f^{-1}(U)\to U\) along \(U'\to U\), hence an isomorphism. As isomorphisms can be checked on geometric stalks, the base change map over \(Y'\) is an isomorphism.

Now \(Y\) is affine, hence quasi-compact and quasi-separated, so Theorem \bpout{A.25} with \(S=Y\) gives a proper surjection \(\pi:\widetilde X\to X\) and closed immersions \(i_Y:\widetilde X\hookrightarrow\mathbf P^n_Y\) and \(i_X=(\pi,i_Y):\widetilde X\hookrightarrow\mathbf P^n_X\) with \(f\pi\) equal to \(i_Y\) followed by the projection. Here we need Chow's lemma over a base that is affine but need not be Noetherian; we prove this form in Appendix \bpout{A}. The two maps \(\pi\) and \(f\pi\) factor as closed immersions \(i_X\) and \(i_Y\) followed by projections of projective spaces. Closed immersions are finite, so they satisfy \(\mathrm{PBC}_\Lambda\) by Lemma \hyperref[bp80]{4.27}, and projections of projective spaces satisfy \(\mathrm{PBC}_\Lambda\) by the hypothesis and Lemma \hyperref[bp91]{4.31}. By stability under composition (Lemma lemma-composition-cancellation), both \(\pi\) and \(f\pi\) therefore satisfy \(\mathrm{PBC}_\Lambda\). As \(\pi\) is surjective, we may cancel it by Lemma lemma-composition-cancellation, and this gives \(\mathrm{PBC}_\Lambda\) for \(f\).
\end{proof}


% object: sec-Pro-etale-complexes (blueprint/chap-Proper-base-change/sec-Pro-etale-complexes/_section.md), content hidden
\subsection*{\texorpdfstring{4.6\quad Pro-étale complexes}{4.6 Pro-étale complexes}}
\phantomsection\label{bp96}%
\addcontentsline{toc}{subsection}{\texorpdfstring{4.6\quad Pro-étale complexes}{4.6 Pro-étale complexes}}
% text: prose-04-35
\phantomsection\label{bp97}%
The six operations live on the pro-étale site, where we also meet torsion sheaves of modules over rings that are not torsion, such as \(\mathcal O_E\). So far \(\mathrm{PBC}_\Lambda(f)\) is known on the étale site only when \(\Lambda\) is torsion (Propositions \hyperref[bp76]{4.26} and \hyperref[bp94]{Proposition 4.32}). We therefore drop the condition on \(\Lambda\), keeping torsion cohomology sheaves, and pass to classical pro-étale complexes by the comparison theorems of Chapter \bpout{2}.

The reduction to torsion rings is possible because a torsion sheaf of \(\Lambda\)-modules is a filtered colimit of sheaves killed by an integer \(n\). Such a sheaf has a finite filtration with quotients killed by primes \(p\), and these quotients are modules over \(\Lambda/p\Lambda\). Moreover, the cohomology of a sheaf of modules is that of its underlying abelian sheaf, so the base change map does not depend on which ring we regard as the coefficients.

For the passage to the pro-étale site we use two facts: a classical complex in \(D^+\) is \(\nu_X^*\) of an étale one (Corollary \bpout{2.37}), and \(\nu^*\) commutes with \(g^*\) and with direct images along proper maps (Theorem \bpout{2.38}). Taking for \(g\) a geometric point then shows that the stalks of \(Rf_*K\) are the cohomology of the geometric fibres (Corollary \hyperref[bp102]{4.34}).

% object: theorem-proper-base-change-classical-pro (blueprint/chap-Proper-base-change/sec-Pro-etale-complexes/theorem-proper-base-change-classical-pro/_section.md)
\phantomsection\label{bp98}%
\begin{bpstatementit}{Theorem~4.33}{ (Pro-étale proper base change)}
Let \(\Lambda\) be a commutative ring, let \(f:X\to Y\) be a proper morphism of schemes, and let \(g:Y'\to Y\) be arbitrary. Set \(X'=X\times_YY'\), with projections \(f':X'\to Y'\) and \(g':X'\to X\). Let \(K\in D^+_{\mathrm{class}}(X_{\mathrm{pro}\acute{e}t},\Lambda)\) have torsion cohomology sheaves (as abelian sheaves). Then the base change map (4.5) of Definition \hyperref[bp18]{4.6}

\[
g^*Rf_*K\longrightarrow Rf'_*g'^*K
\]

is an isomorphism in \(D^+_{\mathrm{class}}(Y'_{\mathrm{pro}\acute{e}t},\Lambda)\).
\end{bpstatementit}

% text: prose-04-36
\phantomsection\label{bp99}%
Suppose \(\Lambda\) is killed by a positive integer, for instance that \(\Lambda\) is a finite ring. Then every complex of \(\Lambda_X\)-modules has torsion cohomology sheaves, and the theorem applies to every \(K\in D^+_{\mathrm{class}}(X_{\mathrm{pro}\acute{e}t},\Lambda)\).

% proof: proof/theorem-proper-base-change-classical-pro
\begin{proof}[{Proof of Theorem 4.33}]\phantomsection\label{bp100}%
We work on the étale site, where Propositions \hyperref[bp76]{4.26} and \hyperref[bp94]{Proposition 4.32} apply to the torsion rings \(\Lambda/p\Lambda\). By dévissage we reach all torsion sheaves of \(\Lambda\)-modules, and the comparison theorems then bring us back to the pro-étale site.

By Corollary \bpout{2.37}, we may write \(K=\nu_X^*K_{\acute{e}t}\) with \(K_{\acute{e}t}\in D^+(X_{\acute{e}t},\Lambda)\), and the cohomology sheaves of \(K_{\acute{e}t}\) are torsion. Indeed, \(\nu_X^*\) is exact, commutes with colimits and is fully faithful on sheaves, and therefore detects the condition \(F=\varinjlim_nF[n]\) of Definition \hyperref[bp36]{4.13}. We must show that the base change map for \(K_{\acute{e}t}\) is an isomorphism in \(D^+(Y'_{\acute{e}t},\Lambda)\), and then pass to the pro-étale site.

We reduce to the rings \(\Lambda/p\Lambda\), \(p\) prime. These are killed by \(p\), so \(\mathrm{PBC}_{\Lambda/p\Lambda}(f)\) holds by Propositions \hyperref[bp76]{4.26} and \hyperref[bp94]{Proposition 4.32}. Now a sheaf \(F\) of \(\Lambda\)-modules with \(pF=0\) is a sheaf of \(\Lambda/p\Lambda\)-modules, and restriction of scalars along \(\Lambda\to\Lambda/p\Lambda\) is exact and conservative. Moreover, it commutes with inverse images, derived direct images and base change maps, since both derived direct images are computed on underlying abelian sheaves (Lemma \hyperref[bp8]{4.2}). Hence the base change map for \(F\) is an isomorphism.

If \(nF=0\), we argue by induction on the number of prime factors of \(n\). For a prime \(p\mid n\), the outer terms of \(0\to F[p]\to F\xrightarrow{p}pF\to0\) are killed by \(p\) and by \(n/p\), so the five lemma applies to the long exact cohomology sequences. Finally, a torsion sheaf is the filtered colimit \(\varinjlim_nF[n]\). Since (4.7) applies to the qcqs morphisms \(f\) and \(f'\), the base change map is an isomorphism for every torsion sheaf of \(\Lambda\)-modules. The truncation argument at the end of the proof of Theorem \hyperref[bp45]{4.17} then passes from sheaves to bounded below complexes with torsion cohomology, in particular to \(K_{\acute{e}t}\).

To return to the pro-étale site we apply \(\nu_{Y'}^*\). Since \(f\) and \(f'\) are proper, hence qcqs, the comparison of direct images in Theorem \bpout{2.38} applies. Using it together with the comparison of inverse images \(g_{\mathrm{pro}\acute{e}t}^*\nu_Y^*\simeq\nu_{Y'}^*g_{\acute{e}t}^*\), we identify \(\nu_{Y'}^*g^*Rf_*K_{\acute{e}t}\) and \(\nu_{Y'}^*Rf'_*g'^*K_{\acute{e}t}\) with \(g^*Rf_*K\) and \(Rf'_*g'^*K\) in \(D^+(Y'_{\mathrm{pro}\acute{e}t},\Lambda)\). By the last assertion of Lemma \hyperref[bp20]{4.7}, under these identifications the base change map of \(K\) is the inverse image on \(Y'_{\mathrm{pro}\acute{e}t}\) of the étale base change map of \(K_{\acute{e}t}\). We have shown that the latter is an isomorphism.
\end{proof}

% text: prose-04-37
\phantomsection\label{bp101}%
Taking for \(g\) a geometric point, we answer the question at the start of the chapter. The only additional point is to make sense of a stalk on the pro-étale site, where the stalk is the value of \(\bar y^*Rf_*K\) on the point \(\operatorname{Spec}\Omega\), whose topos is that of condensed sets. We check that for a classical complex this value is the étale stalk.

% object: corollary-cohomology-geometric-fibres (blueprint/chap-Proper-base-change/sec-Pro-etale-complexes/corollary-cohomology-geometric-fibres/_section.md)
\phantomsection\label{bp102}%
\begin{bpstatementit}{Corollary~4.34}{ (Cohomology of geometric fibres)}
Under the hypotheses of Theorem \hyperref[bp98]{4.33}, for every geometric point \(\bar y:\operatorname{Spec}(\Omega)\to Y\) there is a canonical isomorphism

\[
(Rf_*K)_{\bar y}\simeq
R\Gamma\bigl(X_{\bar y,\mathrm{pro}\acute{e}t},K|_{X_{\bar y}}\bigr).
\]

For a classical torsion sheaf \(F\) of \(\Lambda\)-modules this says

\[
(R^qf_*F)_{\bar y}\simeq
H^q_{\mathrm{pro}\acute{e}t}(X_{\bar y},F|_{X_{\bar y}})\qquad(q\geq0).
\]
\end{bpstatementit}

% proof: proof/corollary-cohomology-geometric-fibres
\begin{proof}[{Proof}]\phantomsection\label{bp103}%
We apply the theorem to \(g=\bar y\) and take global sections over the point. The right side becomes the cohomology of the fibre, and we check that the left side is the étale stalk.

Take \(Y'=\operatorname{Spec}(\Omega)\) in the theorem and apply the functor of global sections on \(Y'_{\mathrm{pro}\acute{e}t}\). By Proposition \bpout{2.50} the topos of \(Y'_{\mathrm{pro}\acute{e}t}\) is that of condensed sets, where global sections is evaluation at the point. This functor is exact, because every surjection of profinite sets onto a point has a section. On the right we therefore get the cohomology of the fibre. On the left, \((Rf_*K)_{\bar y}\) means \(\Gamma(Y'_{\mathrm{pro}\acute{e}t},\bar y^*Rf_*K)\), and by the proof of the theorem we have \(Rf_*K\simeq\nu_Y^*Rf_*K_{\acute{e}t}\). Since \(\bar y^*\nu_Y^*\simeq\nu_{Y'}^*\bar y^*\), the complex \(\bar y^*Rf_*K\) is the constant condensed complex on the étale stalk \((Rf_*K_{\acute{e}t})_{\bar y}\), so its global sections are that stalk. Finally, by the theorem on the étale site, this stalk is the étale cohomology of the fibre.
\end{proof}


% object: sec-Finite-coefficients (blueprint/chap-Proper-base-change/sec-Finite-coefficients/_section.md), content hidden
\subsection*{\texorpdfstring{4.7\quad Finite coefficients}{4.7 Finite coefficients}}
\phantomsection\label{bp104}%
\addcontentsline{toc}{subsection}{\texorpdfstring{4.7\quad Finite coefficients}{4.7 Finite coefficients}}
% text: prose-04-38
\phantomsection\label{bp105}%
Let \(A\) be a finite admissible coefficient ring and \(p:X\to Y\) a proper morphism of separated schemes of finite type over the standing field \(k\). For \(C\in D^b_{\mathrm{ctf}}(X,A)\), proper base change identifies the geometric stalks of \(Rp_*C\) with the cohomology of the fibres. From this identification and the classical finiteness theorem we obtain that \(Rp_*C\) is constructible, with a bound on its Tor amplitude. Both facts are needed for the passage to integral and rational coefficients below, whereas direct sums in the unbounded localizing categories are treated in Section \bpout{5.3}.

% object: proposition-proper-direct-image-finite-coefficients (blueprint/chap-Proper-base-change/sec-Finite-coefficients/proposition-proper-direct-image-finite-coefficients/_section.md)
\phantomsection\label{bp106}%
\begin{bpstatementit}{Proposition~4.35}{ (Proper direct image for finite coefficients)}
Let \(A\) be a finite admissible ring, and let \(p:X\to Y\) be a proper morphism of separated schemes of finite type over \(k\). Then \(Rp_*\) maps \(D^b_{\mathrm{ctf}}(X,A)\) into \(D^b_{\mathrm{ctf}}(Y,A)\). More precisely, if \(C\in D^b_{\mathrm{ctf}}(X,A)\) has Tor amplitude in \([a,b]\) and \(\operatorname{cd}_\ell(X_{\bar y})\leq d\) for all geometric points \(\bar y\) of \(Y\), then \(Rp_*C\) has Tor amplitude in \([a,b+d]\).
\end{bpstatementit}

% proof: proof/proposition-proper-direct-image-finite-coefficients
\begin{proof}[{Proof}]\phantomsection\label{bp107}%
By Corollary \bpout{2.44} the étale and pro-étale categories \(D^b_{\mathrm{ctf}}\) agree, and by Theorem \bpout{2.38} so do the proper direct images of their bounded below complexes. We can therefore transport bounded constructibility and the Tor bound from the étale site.

On the étale site, let \(C\in D^b_{\mathrm{ctf}}(X_{\acute{e}t},A)\). The proper direct image of a constructible torsion sheaf is constructible [Stacks, \href{https://stacks.math.columbia.edu/tag/0GL0}{0GL0}], and \(\operatorname{cd}_\ell(p)<\infty\) [Stacks, \href{https://stacks.math.columbia.edu/tag/0F10}{0F10}]. Hence \(Rp_*C\in\Dbc(Y_{\acute{e}t},A)\).

To check finite Tor amplitude, suppose that \(C\) has Tor amplitude \([a,b]\), and let \(d\) be an integer with \(\operatorname{cd}_\ell(X_{\bar y})\leq d\) for all geometric points \(\bar y\) of \(Y\). Since derived tensor products are computed stalkwise, it suffices to show that each stalk \((Rp_*C)_{\bar y}\) has Tor amplitude in \([a,b+d]\) as a complex of \(A\)-modules. In other words, \((Rp_*C)_{\bar y}\otimes^{\mathbf L}_AN=((Rp_*C)\otimes_A^{\mathbf L}N_Y)_{\bar y}\) should have cohomology in \([a,b+d]\) for every \(A\)-module \(N\). By proper base change and the projection formula for constant modules [Stacks, \href{https://stacks.math.columbia.edu/tag/0F12}{0F12}], this stalk is

\[
R\Gamma\bigl(X_{\bar y},
C|_{X_{\bar y}}\otimes_A^{\mathbf L}N_{X_{\bar y}}\bigr).
\]

The complex \(C|_{X_{\bar y}}\otimes_A^{\mathbf L}N_{X_{\bar y}}\) has cohomology in \([a,b]\), because restriction to the fibre preserves Tor amplitude. Since \(R\Gamma(X_{\bar y},-)\) preserves \(D^{\geq a}\) and has cohomological dimension at most \(d\), it keeps the lower bound and raises the upper bound by at most \(d\). Thus \(Rp_*C\in D^b_{\mathrm{ctf}}(Y_{\acute{e}t},A)\), with Tor amplitude in \([a,b+d]\).
\end{proof}


% object: sec-ell-adic-coefficients (blueprint/chap-Proper-base-change/sec-ell-adic-coefficients/_section.md), content hidden
\subsection*{\texorpdfstring{4.8\quad $\ell$-adic coefficients}{4.8 ell-adic coefficients}}
\phantomsection\label{bp108}%
\addcontentsline{toc}{subsection}{\texorpdfstring{4.8\quad $\ell$-adic coefficients}{4.8 ell-adic coefficients}}
% text: prose-04-39
\phantomsection\label{bp109}%
Pro-étale proper base change (Theorem \hyperref[bp98]{4.33}) concerns classical complexes with torsion cohomology sheaves. For a finite ring \(A\) it therefore applies, with \(\Lambda=A\), to every \(K\in D^b_{\mathrm{ctf}}(X,A)\) for schemes of finite type over \(k\), because bounded constructible complexes on quasi-compact quasi-separated schemes are classical (Theorem \bpout{2.43}). For \(\mathcal O_E\), however, the theorem gives proper base change only at each finite level \(\mathcal O_E/\pi_E^n\), and we must pass to the limit. Later statements of this kind make this passage once and for all, by the reduction principle of Section \bpout{7.3}; here we make it directly.

The difficulty is completeness: the inverse image \(g^*:D(Y_{\mathrm{pro}\acute{e}t},\mathcal O_E)\to D(Y'_{\mathrm{pro}\acute{e}t},\mathcal O_E)\), which is exact on sheaves of modules, need not preserve derived completeness. We therefore complete after pulling back.

% object: definition-completed-inverse-image (blueprint/chap-Proper-base-change/sec-ell-adic-coefficients/definition-completed-inverse-image/_section.md)
\phantomsection\label{bp110}%
\begin{bpstatement}{Definition~4.36}{ (Completed inverse image)}
Let \(g:Y'\to Y\) be a morphism. For \(M\in D(Y_{\mathrm{pro}\acute{e}t},\mathcal O_E)\), the \emph{completed inverse image} is

\[
\hat g^*M
=R\varprojlim_n g^*(M\otimes^{\mathbf L}_{\mathcal O_{E,Y}}
(\mathcal O_E/\pi_E^n)_Y)
\in D(Y'_{\mathrm{pro}\acute{e}t},\mathcal O_E).
\]
\end{bpstatement}

% text: prose-04-40
\phantomsection\label{bp111}%
The natural map \(g^*M\to\hat g^*M\) is the completion map of \(g^*M\), since \(g^*\) commutes with the reductions. For constructible \(M\) the completion does nothing. Indeed, let \(Y\) and \(Y'\) be of finite type over \(k\) and \(M\in D^b_{\mathrm{ctf}}(Y,\mathcal O_E)\). Then \(g^*M\) lies in \(\Dbc(Y',\mathcal O_E)\) by Lemma lemma-integral-models, hence it is derived complete by Proposition \bpout{3.11} ((1)\(\Rightarrow\)(2)), and \(\hat g^*M=g^*M\). However, until we know that proper direct image preserves constructibility, we can only state base change with the completed inverse image.

In the proof we use the following map, which also appears in the projection formula.

% object: definition-projection-formula-map (blueprint/chap-Proper-base-change/sec-ell-adic-coefficients/definition-projection-formula-map/_section.md)
\phantomsection\label{bp112}%
\begin{bpstatement}{Definition~4.37}{ (Projection formula map)}
Let \(p:X\to Y\) be a morphism, \(K\in D(X_{\mathrm{pro}\acute{e}t},\Lambda)\) and \(M\in D(Y_{\mathrm{pro}\acute{e}t},\Lambda)\). The \emph{projection formula map}

\[
Rp_*K\otimes^{\mathbf L}M\longrightarrow Rp_*(K\otimes^{\mathbf L}p^*M)
\tag{4.14}
\]

is the adjoint, under \(p^*\dashv Rp_*\), of the composite \(p^*(Rp_*K\otimes^{\mathbf L}M)\simeq p^*Rp_*K\otimes^{\mathbf L}p^*M
\to K\otimes^{\mathbf L}p^*M\) of the monoidal isomorphism of \(p^*\) and the counit \(p^*Rp_*K\to K\) tensored with \(p^*M\).
\end{bpstatement}

% object: proposition-proper-base-change-complete-integral (blueprint/chap-Proper-base-change/sec-ell-adic-coefficients/proposition-proper-base-change-complete-integral/_section.md)
\phantomsection\label{bp113}%
\begin{bpstatementit}{Proposition~4.38}{ (Proper base change for complete complexes)}
Let \(p:X\to Y\) be proper, let \(g:Y'\to Y\) be any morphism, and let \(p':X'\to Y'\) and \(g':X'\to X\) be the projections from \(X'=X\times_YY'\). Fix a finite extension \(E/\mathbf Q_\ell\). Suppose that \(K\in D(X_{\mathrm{pro}\acute{e}t},\mathcal O_E)\) is derived \(\pi_E\)-complete and that each reduction

\[
K_n=K\otimes^{\mathbf L}_{\mathcal O_{E,X}}
(\mathcal O_E/\pi_E^n)_X
\]

belongs to \(D^+_{\mathrm{class}}(X_{\mathrm{pro}\acute{e}t},\mathcal O_E/\pi_E^n)\). Then \(Rp_*K\) is derived \(\pi_E\)-complete, and there are canonical isomorphisms

\[
\hat g^*Rp_*K\simeq R\varprojlim_n g^*Rp_*K_n,
\qquad
Rp'_*\hat g'^*K\simeq R\varprojlim_n Rp'_*g'^*K_n,
\]

and under them the derived limit of the base change maps \(g^*Rp_*K_n\to Rp'_*g'^*K_n\) is an isomorphism \(\hat g^*Rp_*K\xrightarrow{\ \sim\ }Rp'_*\hat g'^*K\).
\end{bpstatementit}

% proof: proof/proposition-proper-base-change-complete-integral
\begin{proof}[{Proof}]\phantomsection\label{bp114}%
Theorem \hyperref[bp98]{4.33} applies at each finite level, so the point is that \(Rp_*\) and the completed inverse images are compatible with the passage to the derived limit of the reductions.

The reductions of \(Rp_*K\) are controlled by the projection formula. The projection formula map (4.14) for \(M=(\mathcal O_E/\pi_E^n)_Y\) is a map

\[
(Rp_*K)\otimes^{\mathbf L}(\mathcal O_E/\pi_E^n)_Y
\longrightarrow Rp_*K_n,
\]

which we regard in \(D(Y_{\mathrm{pro}\acute{e}t},\mathcal O_E/\pi_E^n)\). We claim that it is an isomorphism. Since restriction of scalars to \(\mathcal O_E\) is conservative, we may check this after forgetting the \(\mathcal O_E/\pi_E^n\)-structure. Now the ring \(\mathcal O_E/\pi_E^n\) has the free resolution \([\mathcal O_E\xrightarrow{\pi_E^n}\mathcal O_E]\), so by Corollary \bpout{2.52} the source is the cone of \(\pi_E^n\) on \(Rp_*K\), and the target is \(Rp_*\) of the cone of \(\pi_E^n\) on \(K\). These agree because \(Rp_*\) is triangulated and \(\mathcal O_E\)-linear. Moreover, the resolutions for \(n+1\) and \(n\) are related by the chain map \((\pi_E,\mathrm{id})\), and applying \(Rp_*\) to it shows that these comparisons form an isomorphism of towers. Since all these constructions can be made on complexes with functorial resolutions, the comparisons define a map in the derived category of towers of Definition \bpout{2.61}.

Since \(K\) is derived complete, it is the derived limit of its reductions \(K_n\), and since the functor \(Rp_*\) is a right adjoint, it commutes with derived limits, and we obtain

\[
Rp_*K\simeq Rp_*R\varprojlim_nK_n
\simeq R\varprojlim_nRp_*K_n.
\]

By the reduction comparisons, the right side is the derived limit of the reductions \((Rp_*K)\otimes^{\mathbf L}(\mathcal O_E/\pi_E^n)_Y\), that is, the derived completion of \(Rp_*K\). Hence the completion map of \(Rp_*K\) is an isomorphism, and \(Rp_*K\) is derived complete. The definition of completed inverse image and the same right-adjoint property for \(Rp'_*\) now give

\[
\begin{aligned}
\hat g^*Rp_*K
&\simeq R\varprojlim_n g^*Rp_*K_n,\\
Rp'_*\hat g'^*K
&\simeq R\varprojlim_n Rp'_*g'^*K_n.
\end{aligned}
\]

The base change maps \(g^*Rp_*K_n\to Rp'_*g'^*K_n\) are natural in \(K_n\), so they form a map of towers, and each is an isomorphism by Theorem \hyperref[bp98]{4.33}. Its cone therefore has zero cohomology at every level, and applying \(R\varprojlim\) gives the asserted isomorphism.
\end{proof}

% text: prose-04-41
\phantomsection\label{bp115}%
When the pullbacks are already complete, the completed base change isomorphism is the ordinary one.

% object: lemma-completed-ordinary-base-change (blueprint/chap-Proper-base-change/sec-ell-adic-coefficients/lemma-completed-ordinary-base-change/_section.md)
\phantomsection\label{bp116}%
\begin{bpstatementit}{Lemma~4.39}{ (Completed and ordinary base change)}
In the setting of Proposition \hyperref[bp113]{4.38}, the square formed by the base change map \(g^*Rp_*K\to Rp'_*g'^*K\) of (4.5), the completed base change isomorphism \(\hat g^*Rp_*K\to Rp'_*\hat g'^*K\), and the completion maps \(g^*Rp_*K\to\hat g^*Rp_*K\) and \(Rp'_*g'^*K\to Rp'_*\hat g'^*K\) commutes. In particular, if \(g^*Rp_*K\) and \(g'^*K\) are derived complete, then the base change map \(g^*Rp_*K\to Rp'_*g'^*K\) is an isomorphism.
\end{bpstatementit}

% proof: proof/lemma-completed-ordinary-base-change
\begin{proof}[{Proof}]\phantomsection\label{bp117}%
Both maps are built from the counit of \(p^*\dashv Rp_*\), as in Definition \hyperref[bp18]{4.6}, and this counit is natural with respect to the quotient maps and the chain maps \((\pi_E,\mathrm{id})\) relating the reductions. Hence the square commutes.

Now suppose that \(g^*Rp_*K\) and \(g'^*K\) are derived complete. Then the first completion map is an isomorphism, and the second is \(Rp'_*\) applied to the completion map of \(g'^*K\), so it is an isomorphism as well. By the commutative square, the base change map is then a composite of isomorphisms.
\end{proof}

% text: prose-04-42
\phantomsection\label{bp118}%
We will show that \(Rp_*K\) is again constructible when \(K\) is constructible. Once this is known, the completed inverse images in Proposition \hyperref[bp113]{4.38} agree with \(g^*\) and \(g'^*\), and the proposition becomes the ordinary base change isomorphism

\[
g^*Rp_*K\xrightarrow{\sim}Rp'_*g'^*K
\qquad(K\in\Dbc(X,\mathcal O_E)).
\]


% object: sec-Proper-direct-image (blueprint/chap-Proper-base-change/sec-Proper-direct-image/_section.md), content hidden
\subsection*{\texorpdfstring{4.9\quad Proper direct image}{4.9 Proper direct image}}
\phantomsection\label{bp119}%
\addcontentsline{toc}{subsection}{\texorpdfstring{4.9\quad Proper direct image}{4.9 Proper direct image}}
% text: prose-04-43
\phantomsection\label{bp120}%
Let \(p:X\to Y\) be proper. We need to know, for every admissible coefficient ring, that \(Rp_*\) preserves \(D^b_{\mathrm{ctf}}\) and satisfies proper base change and the projection formula there, because the construction of \(Rf_!=Rp_*\circ j_!\) in Section \bpout{5.4} starts from these facts. For \(Y=\operatorname{Spec}k\) with \(k\) algebraically closed, they say that the cohomology groups \(H^i(X,\overline{\mathbf Q}_\ell)\) of a proper \(X\) are finite-dimensional and vanish outside a bounded range of degrees.

For finite coefficients, constructibility and the Tor bound are Proposition \hyperref[bp106]{4.35}, and the other admissible rings reduce to this case.

For \(\mathcal O_E\), the complex \(Rp_*K\) is derived complete, and its reduction modulo \(\pi_E\) is \(Rp_*\) of the reduction of \(K\), so the integral constructibility criterion (Proposition \bpout{3.11}) and derived Nakayama (Lemma \bpout{2.62}) reduce everything to torsion coefficients. For \(E\) and \(\overline{\mathbf Q}_\ell\) it suffices to treat the generators \((a_!L_0)[1/\ell]\). Indeed, on bounded below complexes \(Rp_*\) commutes with inverting \(\ell\) and with extension of scalars to finite extensions and to \(\overline{\mathbf Q}_\ell\), since both are filtered colimits (Corollary \bpout{2.29}).

% object: proposition-proper-direct-image-constructible-complexes (blueprint/chap-Proper-base-change/sec-Proper-direct-image/proposition-proper-direct-image-constructible-complexes/_section.md)
\phantomsection\label{bp121}%
\begin{bpstatementit}{Proposition~4.40}{ (Proper direct image of constructible complexes)}
Let \(p:X\to Y\) be proper and let \(\Lambda\) be admissible.

\begin{enumerate}
\item \(Rp_*:D(X_{\mathrm{pro}\acute{e}t},\Lambda)\to D(Y_{\mathrm{pro}\acute{e}t},\Lambda)\) maps \(D^b_{\mathrm{ctf}}(X,\Lambda)\) into \(D^b_{\mathrm{ctf}}(Y,\Lambda)\).
\item For a cartesian square with \(g:Y'\to Y\), \(g':X'\to X\) and \(p':X'\to Y'\), the base change map (4.5) \(g^*Rp_*\to Rp'_*g'^*\) is an isomorphism of functors \(D^b_{\mathrm{ctf}}(X,\Lambda)\to D^b_{\mathrm{ctf}}(Y',\Lambda)\).
\item The projection formula map (4.14) is an isomorphism for \(K\in D^b_{\mathrm{ctf}}(X,\Lambda)\) and \(M\in D^b_{\mathrm{ctf}}(Y,\Lambda)\).
\end{enumerate}
\end{bpstatementit}

% proof: proof/proposition-proper-direct-image-constructible-complexes
\begin{proof}[{Proof}]\phantomsection\label{bp122}%
\emph{Finite coefficients.} Part (1), including the Tor bound, is Proposition \hyperref[bp106]{4.35}, and proper base change holds by Theorem \hyperref[bp98]{4.33} with \(\Lambda=A\). The projection formula, finally, is the classical torsion formula [Stacks, \href{https://stacks.math.columbia.edu/tag/0F0G}{0F0G}] in \(D(Y_{\acute{e}t},A)\), transported to \(D(Y_{\mathrm{pro}\acute{e}t},A)\) by pullback from the étale site.

\emph{Integral coefficients.} Let \(K\in\Dbc(X,\mathcal O_E)\). We first show that \(Rp_*K\in\Dbc(Y,\mathcal O_E)\), by computing its reduction as in (4.15) below. By Proposition \bpout{3.11} and Lemma \bpout{3.5}, \(K\) is derived complete with bounded classical finite reductions. Hence Proposition \hyperref[bp113]{4.38} applies. It shows that \(Rp_*K\) is derived complete and gives a canonical comparison of reductions

\[
\rho_K:(Rp_*K)\otimes^{\mathbf L}\mathcal O_{E,Y}/\pi_E
\xrightarrow{\ \sim\ } Rp_*(K\otimes^{\mathbf L}\mathcal O_{E,X}/\pi_E).
\tag{4.15}
\]

The target of \(\rho_K\) is constructible and bounded. Indeed, it is the direct image of a torsion complex pulled back from \(X_{\acute{e}t}\) [Stacks, \href{https://stacks.math.columbia.edu/tag/09B1}{09B1}], so it is the pullback to \(Y_{\mathrm{pro}\acute{e}t}\) of the direct image in \(D(Y_{\acute{e}t},\mathbf F_E)\) [Stacks, \href{https://stacks.math.columbia.edu/tag/09A6}{09A6}]. The latter is constructible by the finiteness theorem for proper direct image with torsion coefficients [Stacks, \href{https://stacks.math.columbia.edu/tag/0GL0}{0GL0}], and bounded because \(p\) has finite cohomological dimension [Stacks, \href{https://stacks.math.columbia.edu/tag/0F10}{0F10}]. Now Proposition \bpout{3.11} gives \(Rp_*K\in\Dbc(Y,\mathcal O_E)\).

Proper base change holds by Lemma \hyperref[bp116]{4.39}, which applies because \(g^*Rp_*K\) and \(g'^*K\) are derived complete, as inverse image preserves integral constructibility (Lemma lemma-integral-models).

The projection formula is tested modulo \(\pi_E\) in the same way. Let \(C\) be the cone of the projection formula map for \(K\in\Dbc(X,\mathcal O_E)\) and \(M\in\Dbc(Y,\mathcal O_E)\). Its source \(Rp_*K\otimes^{\mathbf L}M\) and target \(Rp_*(K\otimes^{\mathbf L}p^*M)\) are integrally constructible, by Lemma lemma-integral-models and the fact just proved that \(Rp_*\) maps \(\Dbc(X,\mathcal O_E)\) into \(\Dbc(Y,\mathcal O_E)\). As objects of \(\Dbc\) they are derived complete, and hence so is \(C\). Derived Nakayama (Lemma \bpout{2.62}) therefore reduces us to showing \(C/\pi_E=0\). In other words, we must show that the reduction of the projection formula map is an isomorphism.

This reduction is the torsion projection formula map, which is an isomorphism by [Stacks, \href{https://stacks.math.columbia.edu/tag/0F0G}{0F0G}]. To identify the two maps, note that the sources and targets correspond by \(\rho_K\) and associativity of extension of scalars, and the maps agree because both are adjoint to the counit \(p^*Rp_*K\to K\), as is \(\rho_K\).

\emph{Rational coefficients.} Here \(Rp_*\) commutes with the operations that produce the generators of Lemma lemma-integral-models from integral complexes. For \(K\in\Dbc(X,\mathcal O_E)\), the complex \(K[1/\ell]\) is the colimit of \(K\xrightarrow{\ell}K\xrightarrow{\ell}\cdots\), and since this system has a common lower bound, Corollary \bpout{2.29} shows that

\[
Rp_*(K[1/\ell])\simeq(Rp_*K)[1/\ell].
\]

Likewise \(Rp_*\) commutes with extension of scalars to a finite \(E'/E\), which is a finite direct sum by Lemma lemma-inverse-image-coefficient-sheaves. By Corollary \bpout{2.29} again, it also commutes with the filtered colimit \(\QbarX\otimes_{E_X}-=\varinjlim_{E'}E'_X\otimes_{E_X}-\). By the integral case, it follows that \(Rp_*\) maps the generators into \(\Dbc(Y,E)\) and \(\Dbc(Y,\overline{\mathbf Q}_\ell)\). Since \(Rp_*\) is exact and \(\Dbc(Y,E)\) and \(\Dbc(Y,\overline{\mathbf Q}_\ell)\) are triangulated (Corollary \bpout{3.20}), \(Rp_*\) preserves \(\Dbc(-,E)\) and \(\Dbc(-,\overline{\mathbf Q}_\ell)\).

For proper base change and the projection formula it is enough to treat the generators. Indeed, both sides of each map are exact functors in each variable, so the objects for which the map is an isomorphism form a triangulated subcategory. On the generators, the maps are obtained from the integral isomorphisms by inverting \(\ell\) and extending scalars. These operations commute with \(Rp_*\) by the preceding paragraph, and with \(g^*\) and \(\otimes^{\mathbf L}\) because they are given by tensoring with the coefficient sheaves. Moreover, the maps so obtained are the base change and projection formula maps, since these scalar comparisons commute with the counits of \(p^*\dashv Rp_*\).
\end{proof}



\clearpage\phantomsection\hypertarget{bp-outside}{}%
\section*{Outside this export}
\addcontentsline{toc}{section}{Outside this export}
References marked \ensuremath{{}^\ast} are to parts of the blueprint this document leaves out:
\begin{description}
\item[Section 1.3] The function–sheaf dictionary, in 1 Introduction (\bpslug{sec-{}\allowbreak{}The-{}\allowbreak{}function-{}\allowbreak{}sheaf-{}\allowbreak{}dictionary})
\item[Chapter 2] The pro-étale site, in Ⅰ The six operations (\bpslug{chap-{}\allowbreak{}The-{}\allowbreak{}pro-{}\allowbreak{}etale-{}\allowbreak{}site})
\item[Corollary 2.29] Continuity of cohomology, in Ⅰ The six operations (\bpslug{corollary-{}\allowbreak{}continuity-{}\allowbreak{}cohomology})
\item[Corollary 2.37] Bounded-below classical complexes, in Ⅰ The six operations (\bpslug{corollary-{}\allowbreak{}bounded-{}\allowbreak{}below-{}\allowbreak{}classical-{}\allowbreak{}complexes})
\item[Theorem 2.38] Comparison of direct images, in Ⅰ The six operations (\bpslug{theorem-{}\allowbreak{}comparison-{}\allowbreak{}direct-{}\allowbreak{}images})
\item[Lemma 2.39] Compatibilities of the comparison of direct images, in Ⅰ The six operations (\bpslug{lemma-{}\allowbreak{}comparison-{}\allowbreak{}direct-{}\allowbreak{}images-{}\allowbreak{}compatibilities})
\item[Theorem 2.43] Comparison for constructible sheaves, in Ⅰ The six operations (\bpslug{theorem-{}\allowbreak{}comparison-{}\allowbreak{}constructible-{}\allowbreak{}sheaves})
\item[Corollary 2.44] Finite Tor amplitude under comparison, in Ⅰ The six operations (\bpslug{corollary-{}\allowbreak{}finite-{}\allowbreak{}tor-{}\allowbreak{}amplitude-{}\allowbreak{}under-{}\allowbreak{}comparison})
\item[Proposition 2.50] The pro-étale site of a separably closed field, in Ⅰ The six operations (\bpslug{proposition-{}\allowbreak{}pro-{}\allowbreak{}etale-{}\allowbreak{}site-{}\allowbreak{}separably-{}\allowbreak{}closed})
\item[Corollary 2.52] Finite quotients and their transition maps, in Ⅰ The six operations (\bpslug{corollary-{}\allowbreak{}finite-{}\allowbreak{}quotients-{}\allowbreak{}their-{}\allowbreak{}transition-{}\allowbreak{}maps})
\item[Lemma 2.54] Inverse image of the coefficient sheaves, in Ⅰ The six operations (\bpslug{lemma-{}\allowbreak{}inverse-{}\allowbreak{}image-{}\allowbreak{}coefficient-{}\allowbreak{}sheaves})
\item[Definition 2.61] Derived adic completeness, in Ⅰ The six operations (\bpslug{definition-{}\allowbreak{}derived-{}\allowbreak{}adic-{}\allowbreak{}completeness})
\item[Lemma 2.62] Reduction and derived Nakayama, in Ⅰ The six operations (\bpslug{lemma-{}\allowbreak{}reduction-{}\allowbreak{}derived-{}\allowbreak{}nakayama})
\item[Lemma 3.5] Bounds for the finite reductions, in Ⅰ The six operations (\bpslug{lemma-{}\allowbreak{}bounds-{}\allowbreak{}finite-{}\allowbreak{}reductions})
\item[Proposition 3.11] The integral constructibility criterion, in Ⅰ The six operations (\bpslug{proposition-{}\allowbreak{}integral-{}\allowbreak{}constructibility-{}\allowbreak{}criterion})
\item[Corollary 3.20] Generators of $D^b_c$ with field coefficients, in Ⅰ The six operations (\bpslug{corollary-{}\allowbreak{}coefficient-{}\allowbreak{}heart-{}\allowbreak{}b})
\item[Section 5.3] Compactness for finite coefficients, in Ⅰ The six operations (\bpslug{sec-{}\allowbreak{}Compactness-{}\allowbreak{}for-{}\allowbreak{}finite-{}\allowbreak{}coefficients})
\item[Section 5.4] Direct image with compact support, in Ⅰ The six operations (\bpslug{sec-{}\allowbreak{}Direct-{}\allowbreak{}image-{}\allowbreak{}with-{}\allowbreak{}compact-{}\allowbreak{}support})
\item[Proposition 5.26] Construction and independence of $Rf_!$, in Ⅰ The six operations (\bpslug{proposition-{}\allowbreak{}construction-{}\allowbreak{}independence-{}\allowbreak{}rf-{}\allowbreak{}shriek})
\item[Section 7.3] Reduction to finite coefficients, in Ⅰ The six operations (\bpslug{sec-{}\allowbreak{}Reduction-{}\allowbreak{}to-{}\allowbreak{}finite-{}\allowbreak{}coefficients})
\item[Chapter 8] Smooth base change, in Ⅰ The six operations (\bpslug{chap-{}\allowbreak{}Smooth-{}\allowbreak{}base-{}\allowbreak{}change})
\item[Corollary 12.26] Trace functions and direct image with compact support, in Ⅱ Duality and the trace formula (\bpslug{corollary-{}\allowbreak{}trace-{}\allowbreak{}functions-{}\allowbreak{}direct-{}\allowbreak{}image-{}\allowbreak{}compact})
\item[Proposition 12.28] The function–sheaf dictionary, in Ⅱ Duality and the trace formula (\bpslug{proposition-{}\allowbreak{}function-{}\allowbreak{}sheaf-{}\allowbreak{}dictionary})
\item[Appendix A] Nagata compactification and Chow's lemma, in Appendices (\bpslug{chap-{}\allowbreak{}Nagata-{}\allowbreak{}compactification-{}\allowbreak{}and-{}\allowbreak{}Chows-{}\allowbreak{}lemma})
\item[Theorem A.25] Chow's lemma, projective form, in Appendices (\bpslug{theorem-{}\allowbreak{}chow-{}\allowbreak{}s-{}\allowbreak{}lemma-{}\allowbreak{}projective-{}\allowbreak{}form})
\item[Appendix C] The Picard scheme of a curve, in Appendices (\bpslug{chap-{}\allowbreak{}The-{}\allowbreak{}Picard-{}\allowbreak{}scheme-{}\allowbreak{}of-{}\allowbreak{}a-{}\allowbreak{}curve})
\item[Appendix D] Cohomology of torsion sheaves on curves, in Appendices (\bpslug{chap-{}\allowbreak{}Cohomology-{}\allowbreak{}of-{}\allowbreak{}torsion-{}\allowbreak{}sheaves-{}\allowbreak{}on-{}\allowbreak{}curves})
\item[Proposition D.2] Smooth proper curves, torsion prime to the characteristic, in Appendices (\bpslug{proposition-{}\allowbreak{}smooth-{}\allowbreak{}proper-{}\allowbreak{}curves-{}\allowbreak{}torsion-{}\allowbreak{}prime})
\end{description}
\clearpage\phantomsection%
\section*{References}
\addcontentsline{toc}{section}{References}
\bpbibitem{BBD82}{bp124}{A. Beilinson, J. Bernstein, and P. Deligne. \emph{Faisceaux pervers}, volume 100 of \emph{Astérisque}. 1982.}

\bpbibitem{BS15}{bp125}{B. Bhatt and P. Scholze. The pro-étale topology for schemes. \emph{Astérisque}, 369:99–201, 2015. \href{https://people.mpim-bonn.mpg.de/scholze/proetale.pdf}{Author's manuscript}.}

\bpbibitem{Con07}{bp126}{B. Conrad. Deligne's notes on Nagata compactifications. \emph{Journal of the Ramanujan Mathematical Society}, 22:205–257, 2007. \href{https://math.stanford.edu/\~conrad/papers/nagatafinal.pdf}{Author's manuscript}.}

\bpbibitem{Del74}{bp127}{P. Deligne. La conjecture de Weil. I. \emph{Publ. Math. IHÉS}, 43:273–307, 1974.}

\bpbibitem{Del80}{bp128}{P. Deligne. La conjecture de Weil. II. \emph{Publ. Math. IHÉS}, 52:137–252, 1980.}

\bpbibitem{Den22}{bp129}{C. Deninger. There is no “Weil-”cohomology theory with real coefficients for arithmetic curves. \emph{Ann. Sc. Norm. Super. Pisa Cl. Sci.}, 2022. To appear. arXiv:\href{https://arxiv.org/abs/2204.02714}{2204.02714}.}

\bpbibitem{Deu41}{bp130}{M. Deuring. Die Typen der Multiplikatorenringe elliptischer Funktionenkörper. \emph{Abh. Math. Sem. Univ. Hamburg}, 14:197–272, 1941.}

\bpbibitem{Eke90}{bp131}{T. Ekedahl. On the adic formalism. In \emph{The Grothendieck Festschrift, Vol. II}, volume 87 of \emph{Progress in Mathematics}, pages 197–218. Birkhäuser, 1990.}

\bpbibitem{Has36}{bp132}{H. Hasse. Zur Theorie der abstrakten elliptischen Funktionenkörper. I, II, III. \emph{J. Reine Angew. Math.}, 175:55–62, 69–88, 193–208, 1936.}

\bpbibitem{Hum72}{bp133}{J. E. Humphreys. \emph{Introduction to Lie Algebras and Representation Theory}, volume 9 of \emph{Graduate Texts in Mathematics}. Springer, 1972.}

\bpbibitem{Ill94}{bp134}{L. Illusie. Autour du théorème de monodromie locale. \emph{Astérisque}, 223:9–57, 1994. §4. \href{https://www.numdam.org/item/AST_1994__223__9_0/}{Original article}.}

\bpbibitem{Kat88}{bp135}{N. M. Katz. Travaux de Laumon. \emph{Astérisque}, 161–162:105–132, 1988. Séminaire Bourbaki, exp. 691.}

\bpbibitem{KW01}{bp136}{R. Kiehl and R. Weissauer. \emph{Weil conjectures, perverse sheaves and \(\ell\)-adic Fourier transform}, volume 42 of \emph{Ergebnisse der Mathematik}. Springer, 2001.}

\bpbibitem{Lau87}{bp137}{G. Laumon. Transformation de Fourier, constantes d'équations fonctionnelles et conjecture de Weil. \emph{Publ. Math. IHÉS}, 65:131–210, 1987.}

\bpbibitem{Mil86a}{bp138}{J. S. Milne. \href{https://www.jmilne.org/math/xnotes/AVs.pdf}{Abelian Varieties}. In \emph{Arithmetic Geometry}. 1986. Corrected online version.}

\bpbibitem{Mil86b}{bp139}{J. S. Milne. \href{https://www.jmilne.org/math/xnotes/JVs.pdf}{Jacobian Varieties}. In \emph{Arithmetic Geometry}. 1986. Corrected online version.}

\bpbibitem{Ran39}{bp140}{R. A. Rankin. Contributions to the theory of Ramanujan's function \(\tau(n)\) and similar arithmetical functions. II. \emph{Proc. Cambridge Philos. Soc.}, 35:357–372, 1939.}

\bpbibitem{SGA4}{bp123}{M. Artin, A. Grothendieck, and J.-L. Verdier, editors. \emph{Théorie des topos et cohomologie étale des schémas}, volumes 269, 270, 305 of \emph{Lecture Notes in Mathematics}. Springer, 1972–1973. SGA 4.}

\bpbibitem{SGA4½}{bp141}{P. Deligne. \emph{Cohomologie étale}, volume 569 of \emph{Lecture Notes in Mathematics}. Springer, 1977. SGA 4½; \emph{Théorèmes de finitude}, §3, and \emph{Rapport sur la formule des traces}, §3, and \emph{Sommes trigonométriques}, §1.}

\bpbibitem{SGA5}{bp142}{A. Grothendieck, editor. \emph{Cohomologie \(\ell\)-adique et fonctions \(L\)}, volume 589 of \emph{Lecture Notes in Mathematics}. Springer, 1977. SGA 5.}

\bpbibitem{SGA7II}{bp143}{P. Deligne and N. Katz, editors. \emph{Groupes de monodromie en géométrie algébrique}, volume 340 of \emph{Lecture Notes in Mathematics}. Springer, 1973. SGA 7 II; Exposé XIII, (2.1.7).}

\bpbibitem{Stacks}{bp144}{The Stacks Project Authors. \href{https://stacks.math.columbia.edu}{The Stacks Project}. Individual results are cited by their stable tags in the text.}

\bpbibitem{Var05}{bp145}{Y. Varshavsky. Lefschetz–Verdier trace formula and a generalization of a theorem of Fujiwara. 2005. Version 2. arXiv:\href{https://arxiv.org/abs/math/0505564}{math/0505564}. \href{https://arxiv.org/abs/math/0505564v2}{Author's preprint}.}

\bpbibitem{Wei48a}{bp146}{A. Weil. \emph{Sur les courbes algébriques et les variétés qui s'en déduisent}. Hermann, Paris, 1948.}

\bpbibitem{Wei48b}{bp147}{A. Weil. \emph{Variétés abéliennes et courbes algébriques}. Hermann, Paris, 1948.}

\bpbibitem{Wei49}{bp148}{A. Weil. Numbers of solutions of equations in finite fields. \emph{Bull. Amer. Math. Soc.}, 55:497–508, 1949.}

\end{document}
